\documentclass[UTF8, fontset=macnew, a4paper, 12pt]{ctexart}

\usepackage{geometry}
\geometry{top=2.5cm, bottom=2.5cm, left=2.5cm, right=2.5cm}
\usepackage{amsmath, amssymb}
\usepackage{booktabs}
\usepackage{multirow}
\usepackage{tabularx}
\usepackage{graphicx}
\usepackage{caption}
\usepackage{xcolor}
\usepackage{float}
\usepackage{enumitem}
\usepackage{fancyhdr}
\usepackage{setspace}
\usepackage{pifont}
\usepackage[hidelinks]{hyperref}

\graphicspath{{../output/figures/}}
\captionsetup{font=small, labelfont=bf, labelsep=quad}
\setlist{itemsep=2pt, parsep=0pt, topsep=4pt}
\onehalfspacing

% ctexart 預設的圖表標號前綴為簡體「图」，改回正體。
\renewcommand{\figurename}{圖}
\renewcommand{\tablename}{表}

\ctexset{
  section/format = \raggedright\zihao{4}\bfseries,
  section/name   = {},
  section/number = \bignum{\value{section}}、,
  subsection/format = \raggedright\zihao{5}\bfseries,
  subsection/number = \chinese{subsection},
  subsection/aftername = 、,
  subsubsection/format = \raggedright\normalsize\bfseries,
}
% 節次採大寫數字（壹貳參…），與內文的「第壹節」及範例論文的體例一致；
% 小節仍用 \chinese（一二三…），對應內文的「之一」「之二」。
\newcommand{\bignum}[1]{\ifcase#1\or 壹\or 貳\or 參\or 肆\or 伍\or
  陸\or 柒\or 捌\or 玖\or 拾\fi}
\renewcommand{\thesection}{\bignum{\value{section}}}
% 交叉引用子節時補上節次，否則 \ref 只印出「三」而無從判斷屬於哪一節。
% \p@subsection 是 LaTeX 寫入 .aux 時加在 \thesubsection 前的前綴。
\renewcommand{\thesubsection}{\chinese{subsection}}
\makeatletter
\renewcommand{\p@subsection}{\thesection 節之}
\makeatother
\setcounter{secnumdepth}{2}   % subsubsection 不編號（標題已自帶（一）（二））

\pagestyle{fancy}
\fancyhf{}
\fancyhead[L]{}
\fancyhead[R]{\small 進度報告（0928）}
\fancyfoot[C]{\small \thepage}
\renewcommand{\headrulewidth}{0.4pt}

% U+2460~2463（\cn{1}~\cn{4}）不在 xeCJK 認定的 CJK 範圍內，會落到拉丁字型而缺字，
% 故改用 pifont 的圈號。
\newcommand{\cn}[1]{\ding{\number\numexpr171+#1\relax}}
\newcommand{\edag}{e_0^{\dagger}}
\newcommand{\LCe}{\mathrm{LC}e^{\dagger}}
\newtheorem{prop}{Proposition}
\newtheorem{cor}{Corollary}
\newtheorem{lem}{Lemma}
\newcommand{\Ex}{\mathbb{E}}

\begin{document}

\begin{center}
{\zihao{3}\bfseries 小人口 Lee-Carter 參數偏誤的閉式校正}\\[4pt]
{\zihao{5} —— 以估計方程的一致性同時改善 $\alpha_x$ 與 $\beta_x$}\\[6pt]
{\zihao{5} Closed-Form Bias Correction for Lee--Carter Parameters in Small
Populations:}\\[2pt]
{\zihao{5} Improving $\alpha_x$ and $\beta_x$ Simultaneously via
Estimating-Equation Consistency}\\[12pt]
{\zihao{5} （林子立）\quad 指導教授：余清祥\,教授}\\[4pt]
{\zihao{5} 國立政治大學統計學系}\\[4pt]
{\zihao{5} 中華民國 115 年 9 月 28 日}
\end{center}

\vspace{6pt}
\begin{center}
\begin{minipage}{0.92\textwidth}
\small
\textbf{摘要}\par\vspace{4pt}
Lee-Carter 模型的年齡參數在小人口下有方向明確的偏誤，既有研究以模擬量化
此一現象並以修勻或參考人口加以緩解。本文改由 $M$-估計的一般性質推導之，
並據以回答一個既有文獻未處理的問題：年齡水準參數 $\alpha_x$ 與年齡敏感度
參數 $\beta_x$ 的偏誤能否\textbf{同時}降低。
標準 Lee-Carter 是取平方損失的 $M$-估計量，而 $M$-估計的 Fisher 一致性要求
估計方程的期望為零；本文指出該條件在小人口下不成立，且不成立的量有閉式
$b(\mu;c_0)=\Ex[\log\max(D,c_0)]-\log\mu$。據此，$\hat\alpha_x$ 的偏誤恰為
該年齡各格 $b$ 的算術平均，此為恆等式而非漸近近似。以混合模型的角度看，
零格替代使每一格的觀測成為一個兩成分混合，其混合權重 $\pi=e^{-\mu}$
\textbf{為已知}，故此處是污染比例已知的 $\varepsilon$-污染情形——這同時說明
穩健損失為何在此付出不必要的效率代價，並導出中位數失效的解析門檻
$\mu<\ln 2$。以 M-分位數 LASSO 的一般目標函數看，其三項設定（損失形狀、
權重、懲罰）皆不改變估計方程的期望，故偏誤校正與三者正交。
臺灣女性 2001--2024 年資料的模擬顯示，該閉式預測逐年齡偏誤的相關係數為
$1.000$、$1.000$、$0.988$；據此修改估計使 $\alpha$ 偏誤中位數下降 $2.3$ 至 $7.2$ 倍。
至於兩參數能否同時改善，本文的分解給出明確的答案：$\alpha_x$ 的偏誤幾乎
全為確定性、$\beta_x$ 的則九成以上為隨機，故校正若逐格施加會把估計 $\mu$
的噪音注入每一格而使 $\beta_x$ 變差；改為僅校正 $\alpha_x$ 則完全不動
中心化矩陣，兩項改善因而可以疊加。人數五萬時該作法使 $\alpha$ 最大偏誤
由 $0.866$ 降至 $0.044$（優於 Firth 的 $0.068$），配合資訊加權後
$\beta$ 的平方誤差和由 $4.449$ 降至 $1.150$，而零歲平均餘命的偏誤
由 $-1.480$ 歲降至 $+0.007$ 歲。
\par\vspace{6pt}
\textbf{關鍵詞}：Lee-Carter 模型、小人口、$M$-估計、Fisher 一致性、
混合模型、對數轉換偏誤
\end{minipage}
\end{center}

\vspace{10pt}

\section{緒論}

\subsection{研究背景}

平均餘命是衡量地區發展與居民健康的重要指標，可協助政府配置醫療與社會福利
資源，也用於保險業的費率計算。我國官方編製的生命表只涵蓋全國與縣市層級，
並未涉及鄉鎮市區。缺漏的原因不在資料不全，而在樣本規模：368 個鄉鎮市區中
有 217 個的單一性別人數在兩萬以下、299 個少於五萬人。

國內對此已有一系列研究。陳政勳與余清祥（2010）探討臺北市、雲嘉兩縣與澎湖縣
的小區域人口推估；王信忠等（2012）處理小區域死亡率的推估；
王信忠等（2017）編製臺灣原住民的死亡率與生命表，該族群的規模與鄉鎮市區相當；
余清祥等（2021）以年輪變動比處理小區域人口推估；
余清祥等（2025）分析平均餘命與標準化死亡率的關係。
在國際期刊上，Wang et al.（2018）與 Yue et al.（2019）則以模擬系統地刻畫了
Lee-Carter（以下簡稱 LC）模型在小人口下的表現。

這一系列研究共同指出兩個\textbf{性質不同}的問題。第一個是\textbf{震盪}：
人數少時年齡別死亡率的觀測值劇烈跳動，死亡率較低的年齡組甚至整年沒有死亡
人數。第二個是\textbf{有方向的偏誤}：Wang et al.（2018）與余清祥等（2021）
指出 LC 的年齡參數在人數減少時會朝同一方向偏移，且在人數二十萬以下時特別
明顯。兩者必須分開處理，因為\textbf{震盪看得見而偏誤看不見}——
震盪可由重複抽樣或加寬區間察覺，偏誤卻在所有樣本上朝同一方向偏移，
單看估計結果無從發現。

圖~\ref{fig:osc} 以臺灣女性 2001--2024 年資料重現這兩個問題，三格依序
呈現它們在三個層次上的表現。

\begin{figure}[H]
\centering
\includegraphics[width=\textwidth]{figT0_oscillation.png}
\caption{小人口下 LC 估計的震盪與偏誤（臺灣女性，以全國配適為真值）。
(a) 單次實現的年齡別死亡率：小人口的曲線劇烈跳動，而空心點為零死亡格，
其被替代後的位置\textbf{系統性地高於}真值曲線。
(b) $\hat\beta_x$ 的 20 次重複：形狀完全不穩定。
(c) 零歲平均餘命的 100 次重複：點雲全部落在真值（虛線）之下。}
\label{fig:osc}
\end{figure}

圖~\ref{fig:osc}(a) 顯示震盪的來源。人數一萬時 1 至 19 歲的各組幾乎全是
零死亡格（空心點），以 $0.5$ 替代後被擺在真值曲線之上約兩個對數單位；
人數五萬時情況緩解但仍有多格為零。值得注意的是這些點\textbf{不是隨機地
散在真值兩側}，而是一致偏高——震盪與偏誤在此是同一個機制的兩面。

圖~\ref{fig:osc}(b) 顯示震盪傳到參數上的結果。$\hat\beta_x$ 描述各年齡對
時間趨勢的敏感度，其真值（黑線）在幼年與 65 至 74 歲有兩個峰；
但 20 次重複的軌跡完全看不出這個形狀，甚至頻繁出現負值。

圖~\ref{fig:osc}(c) 則顯示兩個問題在標的量上同時出現，而且量級相當。
以配適末年的零歲平均餘命為例，真值為 $84.07$ 歲，而標準 LC 的重複中位數
在人數一萬、五萬、二十萬時分別為 $80.45$、$82.59$、$83.17$ 歲，
即系統性低估 $3.62$、$1.48$、$0.90$ 歲；同時 $5$ 至 $95$ 百分位的全距為
$3.00$、$1.84$、$2.01$ 歲。\textbf{人數五萬時偏誤（$1.48$ 歲）與整個隨機
全距（$1.84$ 歲）幾乎相同}，亦即把估計重複一百次所看到的全部變動，
其幅度不過與一個看不見的系統性偏移相當。考慮到我國縣市之間平均餘命的最大
差距約為 $10$ 歲，$1.48$ 歲的偏誤已足以使一個地區在排序上錯置數個名次。
換言之，若不處理偏誤，鄉鎮市區生命表所呈現的地區差異將有相當部分是估計
方法的產物而非真實的健康差異。

\subsection{小樣本估計的兩條傳統}

震盪與偏誤的分工在統計學上有各自的傳統，而本文的取徑須放在這個脈絡中
才說得清楚。

處理\textbf{偏誤}的傳統是小樣本偏誤減少。最大概似估計量在有限樣本下一般
有 $O(n^{-1})$ 的偏誤，Bartlett（1953）與 Cox and Snell（1968）發展出計算
並扣除該項的一般方法，此即\textbf{事後修正}；Firth（1993）則指出事後修正
有一個前提——最大概似估計值必須存在且有限——並在該前提不成立時改採
\textbf{預先修正}，即在分數函數上加一個 Jeffreys 型的懲罰。
這條線的優點是目標明確（就是把偏誤項算出來）且多半無須調參；
缺點是它需要一個可寫下的概似，而且修正的量隨模型而異，沒有通用公式。

處理\textbf{震盪}的傳統則有兩條。一條是 $M$-估計。Huber（1964）把最大概似
推廣為「某個估計方程的解」，使估計量的設計脫離概似的束縛；
其原始動機是穩健性——讓少數離群觀測無法主導估計——
而 Hampel（1974）的影響函數與 Donoho and Huber（1983）的崩潰點則給出了
「能承受多少污染」的度量。$M$-估計的優點是把估計量的性質化約為估計方程的
期望與變異數兩件事，因而可以在不知道真實分布的情形下給出保證；
\textbf{代價是效率}：為了在未知的污染下仍然可靠，它必須放棄一部分在
模型正確時本可獲得的精度。這個代價是否值得，取決於污染比例是否真的未知。

另一條是正則化。Hoerl and Kennard（1970）的 Ridge 與 Tibshirani（1996）的
LASSO 以在目標函數上加懲罰項的方式，以少量偏誤換取大量變異數的降低，
其正當性可追溯至 Stein（1956）的不可容許性與 James and Stein（1961）的
收縮估計。在小區域估計中，Fay and Herriot（1979）的經驗貝氏收縮是同一想法
的區域層級版本。正則化的優點是效果直接且理論成熟；
\textbf{代價是須指定懲罰強度}，而該強度通常以交叉驗證選出——
這在觀測可以自由切分時沒有問題，但在本文的情境下並非如此
（見第\ref{sec:lit}末一小節）。

\subsection{小樣本 $M$-估計的理論發展與應用}

上述兩條線在近二十年逐漸合流，而合流處正是本文所需的工具。

在理論上，$M$-估計的框架被推廣到三個方向。其一是\textbf{與分位數的合併}：
Breckling and Chambers（1988）把 Huber 的損失乘上一個依殘差正負而異的
非對稱權重，得到由穩健程度與非對稱度兩個參數索引的 M-分位數族，
其兩個極限分別是 Koenker and Bassett（1978）的檢查函數與
Newey and Powell（1987）的非對稱最小平方。
其二是\textbf{與懲罰的合併}：Zoubir et al.（2018，第 3 章）整理了穩健損失
與 Lasso／Elastic Net 的結合並給出迴歸係數與尺度的聯合懲罰式 $M$-估計，
Hu et al.（2021）則處理高維情形下的漸近理論。
其三是\textbf{與相依資料的合併}：Liang and Zeger（1986）的廣義估計方程以
工作相關矩陣處理相依而不指定聯合分布，Merlo et al.（2022）進一步把它推廣到
帶 Huber 損失的 M-分位數。

在應用上，這條線在\textbf{小區域估計}中發展得最完整。
Chambers and Tzavidis（2006）以 M-分位數取代 Fay--Herriot 的隨機效果假設，
理由是前者可直接給出各區域自己的係數而不必假設區域效果的分布；
Bianchi et al.（2018）補上推論理論，Ranalli et al.（2023）加入 Lasso 與
Elastic Net 懲罰並應用於空氣品質，Dawber et al.（2022）處理離散標的，
Bugallo et al.（2026）則針對均方誤差估計提出偏誤校正。
Lahiri and Salvati（2023）把巢式誤差迴歸推廣到同時容許區域別係數與區域別
變異成分的情形。這些工作的應用場景多為家戶調查的地方層級指標估計，
其共同特徵是\textbf{以較弱的分布假設換取穩健性}。

\subsection{既有方法在本問題上的適用性}

把上述優缺點對照本文的情境，可以得到三項判斷，而這三項判斷決定了本文的
設計。

\textbf{其一，$M$-估計的效率代價在此不必付。} 穩健方法之所以犧牲效率，
是因為污染比例 $\varepsilon$ 未知，必須為最壞情形投保。
但本文第\ref{sec:method}將指出，小人口 LC 的「污染」來自零格替代這個
\textbf{分析步驟}，其比例恰為卜瓦松的零機率 $e^{-\mu}$，拿到資料即可估計。
既然已知，就不必投保，而可以直接把它造成的後果算出來扣掉。

\textbf{其二，正則化的調參在此無法進行。} 懲罰強度通常以交叉驗證選出，
但 LC 的識別完全依賴年齡與年份兩個方向的跨格借力，切分年齡--年份格會破壞
這個結構，因而沒有合法的切分方式。這使正則化一脈的方法在本問題上
失去其標準的選參程序。

\textbf{其三，$M$-估計真正可借用的不是穩健性，而是更基本的一項性質。}
$M$-估計理論要求估計方程在真值處的期望為零，即 Fisher 一致性。
這個條件在標準教科書中通常一行帶過，因為在其標準設定下它自動成立；
但它自動成立的前提是損失函數作用在\textbf{未經轉換}的反應變數上。
標準 LC 先取對數又替代零格，兩者都是非線性運算，該條件因而不再自動成立。
這正是本文的切入點：把既有研究以模擬觀察到的偏誤，
放回 $M$-估計理論中它本來的位置。

\subsection{研究設計}

基於上述，本文設計一個 LC 的修正估計，其推導依序借用三篇脈絡的性質。

由 \textbf{Huber（1964）} 的 $M$-估計框架，本文把標準 LC 辨識為取平方損失的
$M$-估計量，並依 Fisher 一致性的要求計算其估計方程在真值處的期望；
該期望不為零，其量有閉式，而 $M$-估計對此的標準補救是中心化估計方程。
所提的修正因此不是另行發明，而是這個標準補救在本問題的具體形式。

由\textbf{混合模型}的表示法（McLachlan and Peel 2000；Yao and Xiang 2024），
本文說明零格替代在統計上是什麼：它使每一格成為一個兩成分混合，
而混合權重已知。此一表示同時給出中位數類估計量失效的解析門檻，
並釐清本文的情形與 Lambert（1992）的零膨脹模型有何不同。

由 \textbf{M-分位數與懲罰}的一般目標函數（Breckling and Chambers 1988；
Zoubir et al. 2018；Ranalli et al. 2023），本文說明標準 LC 是該式的一個角落，
而小人口的各種既有補救各自改動了其中一項設定；
據此可判斷本文的修正與它們的關係是正交的，並導出兩項可事先檢驗的預測。

在死亡率文獻中，與本文最接近的是 \textbf{Yue et al.（2019）}，該文以修勻
修改 LC 在小人口的估計。兩者的差別在於：修勻降低的是觀測的變異，
屬處理震盪的一路，且須指定平滑參數；本文的修正處理的是估計方程的中心，
屬處理偏誤的一路，且不引入任何調校參數。
兩者因此互補而非替代，這一點將在第\ref{sec:disc}回到。

\section{文獻回顧}\label{sec:lit}

本節回顧三條脈絡。三者並非死亡率文獻，而是本文推導所依的一般統計理論；
回顧的重點因此放在\textbf{各條脈絡的核心性質}，以及\textbf{該性質在其原始
應用場景與本文情境之間的落差}。第\ref{sec:method}再由這些性質往下推到 LC。

\subsection{$M$-估計與 Fisher 一致性}

Huber（1964）提出 $M$-估計量，把最大概似的想法推廣為「某個目標函數的
極小值」或等價地「某個估計方程的解」：$\hat\theta$ 解
$\sum_i\psi(y_i;\theta)=0$，其中 $\psi$ 不必是對數概似的導數。
該文的動機是位置參數的穩健估計，設定為 $\varepsilon$-污染——絕大多數觀測
來自模型、少數來自別處——目標是使少數觀測無法主導估計。
Hampel（1974）以影響函數刻畫單一觀測的影響力，
Donoho and Huber（1983）以崩潰點刻畫可承受的污染比例上限，
Huber and Ronchetti（2009）為該領域的標準專著。
在應用上，$M$-估計已是訊號處理、電力系統狀態估計與電腦視覺的常規工具
（Zoubir et al. 2018），其場景的共同特徵是觀測數遠大於參數數、
且離群值確實是少數。

$M$-估計的另一條發展線把「估計方程」本身當成建模的對象，而不必先寫下
概似。Godambe 的估計函數理論與 Liang and Zeger（1986）的廣義估計方程
即屬此類，後者以工作相關矩陣處理相依資料而不指定聯合分布，
在縱貫資料分析中已成常規；Merlo et al.（2022）進一步把廣義估計方程推廣到
帶 Huber 損失的 M-分位數。這條線對本文的意義是它確立了一個看法：
一個估計量的性質，可以完全由其估計方程的期望與變異數讀出，
而不必訴諸該估計量是否為某個概似的最大值。本文正是循此讀 LC。

本文所需的性質不是穩健性，而是 $M$-估計理論中更基本的一項：
\textbf{Fisher 一致性}。一個 $M$-估計量要在真值 $\theta_0$ 處一致，
必要條件是估計方程在真值處的期望為零，
$\Ex_{\theta_0}[\psi(Y;\theta_0)]=0$；否則 $\hat\theta$ 收斂到的是
使該期望為零的另一個點，與 $\theta_0$ 相差一個不隨樣本數消失的量。
此一條件在教科書中通常一行帶過，因為對最大概似而言它自動成立
（分數函數的期望為零），對最小平方而言只要誤差期望為零即成立。
但它之所以自動成立，前提是\textbf{損失函數作用在未經轉換的反應變數上}。
一旦反應變數先經過非線性轉換，$\Ex[\psi]=0$ 就不再是自動的——
這正是本文的切入點。

與本文情境的落差因此有三。其一，古典 $M$-估計文獻關心的是
$\psi$ 有界所帶來的穩健性，而本文關心的是 $\psi$ 的\textbf{期望}是否為零，
兩者是不同的問題：平方損失的 $\psi(u)=u$ 完全不穩健，
但若 $\Ex[u]=0$ 它仍是一致的；反之一個很穩健的 $\psi$ 若 $\Ex[\psi]\ne0$
則仍不一致。古典文獻少談後者，因為在其標準設定（損失作用在原始反應變數上）
中它自動成立。其二，古典設定的觀測數遠大於參數數，而 LC 模型每個
$\alpha_x$ 只有 $T$ 個觀測（本文 $T=24$），每個 $\kappa_t$ 只有 $A$ 個
（本文 $A=22$）。其三，古典設定的污染比例未知，而本文的情形不然
（見下一小節）。

還有一項結構上的落差值得先點出。古典 $M$-估計的參數進入模型的方式是線性的
（$x_i^{\top}\theta$），而 LC 的 $\alpha_x+\beta_x\kappa_t$ 對
$(\beta,\kappa)$ 是\textbf{雙線性}的，識別須靠式~\eqref{eq:lc} 的兩個約束。
這使 $\hat\beta$ 與 $\hat\kappa$ 的漸近理論複雜得多。
所幸本文所處理的 $\hat\alpha_x$ 不受此影響——它在雙線性部分被估計之前
即已由列平均定出，因此可以完全在單變數 $M$-估計的框架內分析。
這也是本文把範圍限於 $\alpha_x$ 的技術理由。

\subsection{混合模型與污染比例}

有限混合模型把觀測的分布寫成若干成分分布的加權和，
$f(y)=\sum_{j}\pi_j f_j(y)$。其發展脈絡自 Pearson（1894）以動差法分解
螃蟹前額寬的雙峰分布起，長期受限於計算；
Dempster et al.（1977）的 EM 演算法使最大概似估計成為可行，
此後混合模型成為模型式分群的標準工具，McLachlan and Peel（2000）與
Yao and Xiang（2024）為前後兩代專著，後者涵蓋半參數推廣、穩健建模、
標籤交換與高維情形。在應用上，混合模型用於天文、基因體、經濟與金融等
領域的異質資料；在小區域估計中，Del Sarto et al.（2019）以 M-分位數迴歸的
有限混合處理未觀測異質性，De Nicol\`o et al.（2023）以 Beta 混合估計地區的
不均度指標。

與計數資料最相關的一支是零膨脹模型。Lambert（1992）的零膨脹卜瓦松迴歸把
觀測寫成「一個零處的點質量」與「一個卜瓦松」的兩成分混合，用以處理
製造業瑕疵計數中零的個數遠多於卜瓦松所能解釋者。其形式與本文
第\ref{sec:method}將導出的式~\eqref{eq:mix} 極為相似，因此須把差別講清楚，
否則容易誤讀。零膨脹模型所處理的是\textbf{結構上多出來的零}——
亦即資料生成機制本身含有一個「必然為零」的狀態，故零的個數超過卜瓦松
的預期。本文的情形\textbf{不是}零膨脹：死亡數就是純粹的卜瓦松，
零的個數恰如 $e^{-\mu}$ 所預期，沒有任何多餘。本文的混合是
\textbf{替代規則在取對數之後的變數上造出來的}——因為 $\log 0=-\infty$
不可用，分析者必須把所有零格對應到同一個值 $\log(c_0/E)$，
於是轉換後的變數才出現一個點質量。換言之，零膨脹模型的混合來自
\textbf{資料}，本文的混合來自\textbf{分析步驟}。這個差別有實質後果：
零膨脹模型須估計膨脹比例，而本文的混合權重由卜瓦松本身給定，
無須估計；反之零膨脹模型可以不取對數而直接配適，本文的困難則完全來自
取對數這個步驟。

混合模型與穩健統計的連結早已建立：Huber 的 $\varepsilon$-污染模型
$(1-\varepsilon)F+\varepsilon H$ 就是一個兩成分混合。但兩個文獻的處理方式
相反。混合模型把 $\pi_j$ 與 $f_j$ 都當成待估的參數，以 EM 求解；
穩健統計則把 $\varepsilon$ 與 $H$ 當成未知的滋擾，設計對其不敏感的估計量。
\textbf{穩健估計量之所以必須付出效率代價，正是因為 $\varepsilon$ 未知}：
它必須為最壞情形投保。

與本文情境的落差正在此處，而這個落差對本文有利。
第\ref{sec:method}將指出，零格替代使每一格的觀測成為一個兩成分混合，
而其混合權重是\textbf{已知的}——它就是卜瓦松的零機率 $e^{-\mu}$，
拿到資料即可估計。既然如此，既不必以 EM 去估它，也不必為不知道它而
付保險費。另一項落差則不利：該混合的其中一個成分是\textbf{退化的}
（一個點質量），而有限混合理論通常假設各成分有密度，
故識別性、標籤交換等標準議題在此不適用，EM 亦無從施加。

\subsection{M-分位數與懲罰式估計}

第三條脈絡把損失函數、權重與懲罰整合為一個目標函數。
Koenker and Bassett（1978）以檢查函數 $\rho_\tau$ 取代平方損失，
使迴歸得以刻畫條件分位數；Koenker（2005）為標準參考，
其應用場景包括勞動經濟學的薪資分布與兒童生長參考曲線
（Wei et al. 2006），共同特徵是每個共變數水準上有大量觀測。
Breckling and Chambers（1988）指出，若把 Huber 損失乘上一個依殘差正負
而異的非對稱權重，即得一個由 $(\tau,c)$ 索引的損失函數族 M-分位數，
其兩個極限分別是檢查函數（$c\to0$）與非對稱最小平方
（$c\to\infty$，Newey and Powell 1987）。
Chambers and Tzavidis（2006）把 M-分位數帶進小區域估計，
Bianchi et al.（2018）補上推論理論。

M-分位數之所以在小區域估計中受青睞，理由值得一提，因為它與本文的
取捨相關。傳統的小區域估計以 Fay and Herriot（1979）的隨機效果模型為主，
須假設區域效果的分布（通常為常態）；Chambers and Tzavidis（2006）指出
M-分位數可直接給出各區域自己的迴歸係數而不必作此假設，
因而對分布誤設較不敏感。其應用多為家戶調查資料的地方層級所得與剝奪指標
估計，Bianchi et al.（2018）另處理了推論，Bugallo et al.（2026）則針對
均方誤差估計提出偏誤校正。這一脈的共同特徵是\textbf{以較弱的分布假設
換取穩健性}。本文的取向相反：本文明確採用卜瓦松假設，
並以此換取一個可完全算出的偏誤。兩者各有代價——若卜瓦松假設不成立，
本文的校正就失去依據（見第\ref{sec:disc}）；而若分布假設其實成立，
放棄它就等於放棄可用的資訊。

再加上懲罰即得本文所需的完整形式。Zoubir et al.（2018，第 3 章）
整理了穩健損失與 Lasso／Elastic Net 的結合，包含迴歸係數與尺度的聯合
懲罰式 $M$-估計；Ranalli et al.（2023）以 Lasso 與 Elastic Net 懲罰的
M-分位數迴歸分析交通與氣象對空氣品質的影響；
Hu et al.（2021）另以損失的冪次為刻度提出 $L_p$-分位數並處理高維懲罰。
綜合之，這條脈絡的一般目標函數為
\begin{equation}\label{eq:general}
\min_{\theta}\ \sum_{i} w_i\,\rho_{\tau,c}(u_i(\theta))\;+\;\lambda P(\theta),
\end{equation}
其中含三項設定：損失形狀 $(\tau,c)$、權重 $w$、懲罰 $(\lambda,P)$。

與本文情境的落差有二。其一，這條脈絡的應用場景是單位層級的調查資料，
每個小區域有多個受訪單位；而死亡率資料是聚合的計數，每格只有一個觀測。
其二，三項設定都需要調參（$\tau$、$c$、$\lambda$），
而本問題既無真值、亦無法以交叉驗證選參——切分年齡--年份格會破壞
雙線性結構所依賴的跨格借力。本文因此不使用這條脈絡的方法，
而是借用它的\textbf{分類}：式~\eqref{eq:general} 把標準 LC 納為一個角落
（$\tau=0.5$、$c=\infty$、$w\equiv1$、$\lambda=0$），
因而可以用來判斷各種補救各自動了哪一個設定，以及本文的校正動的是什麼。

\subsection{小結}

表~\ref{tab:three} 把三條脈絡的核心性質、原始應用場景、以及各自對 LC 的
推論並列。第\ref{sec:method}即依此表由上往下推。

\begin{table}[H]
\centering
\caption{三條脈絡的核心性質與對 LC 的推論}
\footnotesize
\renewcommand{\arraystretch}{1.25}
\begin{tabular}{p{2.3cm}p{3.4cm}p{3.5cm}p{4.3cm}}
\toprule
脈絡 & 所借用的核心性質 & 原始應用場景 & 對 LC 的推論\\
\midrule
$M$-估計
& Fisher 一致性：$\Ex[\psi]=0$ 為一致的必要條件（Huber 1964）
& 位置與迴歸的穩健估計；訊號處理、電力系統（Zoubir et al. 2018）
& 標準 LC 是取平方損失的 $M$-估計量，其 $\Ex[\psi]=b(\mu;c_0)\ne0$，
  故\textbf{Fisher 不一致}；標準補救為中心化 $\psi$，即扣除 $b$\\
混合模型
& $\varepsilon$-污染即兩成分混合；崩潰點 $50\%$（Donoho and Huber 1983）
& 模型式分群；零膨脹計數（Lambert 1992）；小區域的異質性（Del Sarto et al. 2019）
& 零格替代使每格成為混合權重\textbf{已知}（$\pi=e^{-\mu}$）的兩成分混合；
  故不必付穩健估計的效率代價，且中位數失效門檻為 $\mu<\ln2$\\
M-分位數\\與懲罰
& 一個含三項設定的目標函數（Breckling and Chambers 1988；Ranalli et al. 2023）
& 家戶調查的小區域指標；高維變數選擇；空氣品質
& 標準 LC 是該式的角落；三項設定（損失函數的形狀、各格的權重、懲罰項）皆不改變
  $\Ex[\psi]$，故偏誤校正與三者\textbf{正交}\\
\bottomrule
\end{tabular}
\label{tab:three}
\end{table}

\subsection{既有方法的限制與本文的檢驗清單}\label{sec:limtable}

上述三條脈絡的每一項工具，其原始文獻都同時記載了優點與限制。
本文的實證部分並非另起爐灶地比較各種方法的優劣，而是\textbf{逐項檢驗這些
已被記載的限制在小人口 LC 上是否確實出現、以及出現時的量級}。
表~\ref{tab:limits} 把這些限制與本文的檢驗位置並列，
後續各節的每一項結果都可回到此表的某一列；第\ref{sec:disc}將逐列回顧。

\begin{table}[H]
\centering
\caption{各方法在原始文獻中所載的限制，以及本文檢驗的位置}
\footnotesize
\renewcommand{\arraystretch}{1.3}
\begin{tabular}{p{2.6cm}p{3.5cm}p{3.6cm}p{4.0cm}}
\toprule
方法（文獻） & 文獻所載的優點 & 文獻所載的限制 & 本文的檢驗與結果\\
\midrule
事後修正\newline（Cox and Snell 1968）
& 目標明確，直接扣除 $O(n^{-1})$ 偏誤項
& \textbf{須最大概似估計值存在且有限}（Firth 1993 明言此點）
& 第\ref{sec:result}：卜瓦松 MLE 在 $N=10^4$ 時為 $-13.03$，
  前提不成立；本文改由 $D$ 的邊際分布推導，繞過此限制\\
預先修正\newline（Firth 1993）
& 最大概似不存在時仍可用
& 為 $O(n^{-1})$ 的漸近結果，且須寫得出概似
& 第\ref{sec:result}：$N=10^4$ 時 $\alpha$ 最大偏誤仍有 $0.427$（非零），
  但遠優於未修正的 $13.03$\\
穩健損失\newline（Huber 1964）
& 不需知道污染的分布即可提供保障
& \textbf{效率代價}（中位數在常態下的相對效率僅 $2/\pi\approx0.64$，
  Huber and Ronchetti 2009）；且假設污染為\textbf{少數}
& 第\ref{sec:method}：污染比例 $\pi=e^{-\mu}$ 已知，故效率代價不必付；
  而幼年組的 $\pi$ 達 $0.74$，「少數」的前提不成立\\
崩潰點\newline（Donoho and Huber 1983）
& 給出方法適用範圍的上界
& 中位數的上界為 $50\%$
& 第\ref{sec:method}：代入混合表示得解析門檻 $\mu<\ln2=0.693$；
  $\mu=0.3$ 時 $\pi=0.741$，與實測的 $74\%$ 零格率相符\\
正則化\newline（Hoerl and Kennard 1970；Tibshirani 1996）
& 以少量偏誤換取大量變異數的降低
& \textbf{須指定懲罰強度}，通常以交叉驗證選出
& 第\ref{sec:lit}：LC 的識別依賴跨格借力，切分年齡--年份格即破壞之；
  第\ref{sec:result}：加權版使 $\beta$ 帶寬收窄七倍卻使中線偏離，
  正是該取捨的具體形式\\
經驗貝氏收縮\newline（Fay and Herriot 1979）
& 收縮量由變異成分定，無須調參
& \textbf{假設直接估計量不偏}，僅有抽樣噪音
& 第\ref{sec:betadec}：$\hat\beta_x$ 本身即有偏（確定性成分佔
  $9.2\%$），前提不完全成立\\
混合模型\newline（McLachlan and Peel 2000）
& 混合權重可由 EM 估計
& 各成分須有密度；EM 另有識別性與標籤交換問題
& 第\ref{sec:method}：本文的成分之一為\textbf{點質量}，EM 無從施加；
  但 $\pi$ 由卜瓦松給定，本就不需估計\\
修勻\newline（Yue et al. 2019）
& 直接降低觀測的震盪
& 須指定平滑參數；處理的是變異數而非偏誤
& 第\ref{sec:disc}：本文的修正不引入平滑參數且處理的是估計方程的中心，
  兩者互補\\
\bottomrule
\end{tabular}
\label{tab:limits}
\end{table}

\subsection{本文後續的安排}\label{sec:roadmap}

依表~\ref{tab:limits} 的檢驗清單，本文後續分四步進行，每一步回答一個
依序遞進的問題。

第\ref{sec:method}處理\textbf{方法論與推導}，回答「標準 LC 違反了哪一個
條件、違反的量是否有閉式」。該節先把標準 LC 辨識為取平方損失的 $M$-估計量，
據 Fisher 一致性的要求算出其估計方程在真值處的期望，得到閉式
$b(\mu;c_0)$，並依 $M$-估計的標準補救導出修正的估計量；
再以混合模型的表示法說明零格替代在統計上是什麼，
並由崩潰點理論導出中位數失效的解析門檻；最後以 M-分位數 LASSO 的一般
目標函數說明該修正與既有補救的關係。該節同時把 $\alpha_x$ 與 $\beta_x$
的偏誤分解為確定性與隨機兩個成分，這個分解決定了後續每一項結果的解讀方式。

第\ref{sec:valid}處理\textbf{模擬驗證}，回答「閉式能否預測實測」。
以全國配適為真值產生小人口，比較閉式預測與模擬所測得的逐年齡偏誤。
此步驟的目的不是比較方法優劣，而是確認推導本身正確；
若閉式與實測不符，後續的修正即失去依據。

第\ref{sec:result}處理\textbf{實證分析}，回答「修正在臺灣資料上能改善多少、
以及改善不到的部分是什麼」。除以卜瓦松模擬比較各估計量外，
另以實際觀測死亡數的二項抽薄產生小人口，藉以檢驗結論在模型誤設下是否穩固。
該節的每一項結果都對應表~\ref{tab:limits} 的某一列。

第\ref{sec:disc}為\textbf{結論、限制與後續工作}，逐列回顧
表~\ref{tab:limits}，說明哪些文獻所載的限制在本問題上確實出現、
哪些因情境不同而得以繞過，並列出本文推導尚未涵蓋的部分。

\section{研究方法}\label{sec:method}

\subsection{模型、記號與標準 LC 的估計}

在推導之前須先把記號與標準作法寫定，因為本文全部的論證都建立在
「標準 LC 的 $\hat\alpha_x$ 究竟是什麼」這個看似瑣碎的細節上。
令 $x=1,\dots,A$ 為年齡組、$t=1,\dots,T$ 為年份，LC 模型設定
\begin{equation}\label{eq:lc}
\log m_{x,t}=\alpha_x+\beta_x\kappa_t,\qquad
\sum_x\beta_x=1,\quad \sum_t\kappa_t=0 .
\end{equation}
其中 $m_{x,t}$ 為中央死亡率，兩個約束用以固定參數的尺度與位置。
觀測為死亡數 $D_{x,t}$ 與曝露數 $E_{x,t}$，而本文設
\begin{equation}\label{eq:pois}
D_{x,t}\sim\text{Poisson}(\mu_{x,t}),\qquad \mu_{x,t}=E_{x,t}\,m_{x,t},
\end{equation}
各格獨立。$D_{x,t}=0$ 時 $\log(0/E)=-\infty$，故實務上必以某正值 $c_0$
替代，慣例 $c_0=0.5$。令
\begin{equation}\label{eq:ell}
\ell_{x,t}=\log\frac{\max(D_{x,t},\,c_0)}{E_{x,t}},\qquad
u_{x,t}(\theta)=\ell_{x,t}-\alpha_x-\beta_x\kappa_t,
\end{equation}
則標準 LC 即式~\eqref{eq:general} 在
$\tau=0.5$、$c=\infty$、$w\equiv1$、$\lambda=0$ 的角落：
\begin{equation}\label{eq:lcobj}
\min_{\theta}\ \sum_{x,t}\tfrac12 u_{x,t}(\theta)^2 ,
\end{equation}
其解由奇異值分解給出，且
\begin{equation}\label{eq:svd}
\hat\alpha_x=\frac1T\sum_t \ell_{x,t},\qquad
(\hat{\boldsymbol\beta},\hat{\boldsymbol\kappa})\ \text{取}\
\{\ell_{x,t}-\hat\alpha_x\}\ \text{的第一組奇異向量} .
\end{equation}

\subsection{估計方程的一致性檢驗}

推導分三步：先寫出 $M$-估計的一般條件，再看標準 LC 落在哪裡，
最後解出估計量實際收斂到的位置。考慮一般的 $M$-估計量，設殘差為 $u_i(\theta)$、損失為 $\rho$、
$\psi=\rho'$，估計量 $\hat\theta$ 使目標函數 $\sum_i\rho(u_i(\theta))$ 極小，
故滿足一階條件
\begin{equation}\label{eq:ee}
\sum_i \psi\big(u_i(\hat\theta)\big)\,
\frac{\partial u_i}{\partial\theta}\Big|_{\hat\theta}=0 .
\end{equation}
在通常的正則條件下，式~\eqref{eq:ee} 除以樣本數後依大數法則收斂到其期望，
故 $\hat\theta$ 收斂到\textbf{母體估計方程}的解 $\theta^{\ast}$，即滿足
\begin{equation}\label{eq:pop}
\Ex\!\left[\psi\big(u(\theta^{\ast})\big)\,
\frac{\partial u}{\partial\theta}\right]=0
\end{equation}
的那個點。所謂 Fisher 一致性，就是要求 $\theta^{\ast}$ 恰為真值 $\theta_0$；
由式~\eqref{eq:pop}，其充要條件是母體估計方程在真值處為零。
此一條件對最大概似自動成立（分數函數的期望為零），
對最小平方而言則只需誤差的期望為零即可——但\textbf{這個「只需」有一個
前提，即損失作用在未經轉換的反應變數上}。一旦反應變數先經過非線性轉換，
誤差的期望是否為零就必須另行計算。

標準 LC 正落在這個「必須另行計算」的情形。由式~\eqref{eq:lcobj}，
平方損失給出 $\psi(u)=u$，而對 $\alpha_x$ 的偏導數為
$\partial u_{x,t}/\partial\alpha_x=-1$，故式~\eqref{eq:pop} 在 $\alpha_x$
這一個分量上化為
\begin{equation}\label{eq:popalpha}
\sum_{t}\Ex\big[\ell_{x,t}-\alpha_x^{\ast}-\beta_x\kappa_t\big]=0 .
\end{equation}
把 $\Ex[\ell_{x,t}]$ 寫成真值加上一個待定的差，即得下述結果。

\begin{prop}\label{prop:fisher}
在式~\eqref{eq:lc} 與~\eqref{eq:pois} 之下，標準 LC 的估計方程在真值處的
期望為
\begin{equation}\label{eq:b}
\Ex\big[\psi(u_{x,t})\big]=\Ex[\ell_{x,t}]-\log m_{x,t}
= b(\mu_{x,t};c_0),
\end{equation}
其中
\begin{equation}\label{eq:bform}
b(\mu;c_0)\;\equiv\;\Ex\big[\log\max(D,c_0)\big]-\log\mu
\;=\;e^{-\mu}\log c_0+\sum_{d\ge1}\frac{e^{-\mu}\mu^{d}}{d!}\log d-\log\mu .
\end{equation}
\end{prop}

證明只需注意曝露數在相減時約掉，
$\ell_{x,t}-\log m_{x,t}=\log\max(D_{x,t},c_0)-\log\mu_{x,t}$，
再按 $D$ 的取值展開期望值：$D=0$ 的機率為 $e^{-\mu}$ 而其貢獻為
$\log c_0$，$D=d\ge1$ 的機率為 $e^{-\mu}\mu^{d}/d!$ 而其貢獻為 $\log d$，
兩者相加再減去 $\log\mu$ 即得式~\eqref{eq:bform}。級數絕對收斂，
實作上取 $d\le\mu+12\sqrt{\mu}+12$ 即達機器精度。

把Proposition~\ref{prop:fisher} 代回式~\eqref{eq:popalpha}，即可解出估計量
實際收斂到的位置：
\begin{equation}\label{eq:astar}
\alpha_x^{\ast}
=\alpha_x+\frac1T\sum_t b(\mu_{x,t};c_0)
\;\equiv\;\alpha_x+\bar b_x .
\end{equation}
式~\eqref{eq:astar} 有一項值得特別強調的後果。$\bar b_x$ 不含樣本數 $T$，
因此\textbf{增加年份並不會使這個偏移消失}——即使觀測一百年，
$\hat\alpha_x$ 仍收斂到 $\alpha_x+\bar b_x$ 而非 $\alpha_x$。
這不是一個隨樣本增大而衰減的小樣本偏誤，而是一個\textbf{不一致}；
能使它縮小的只有人口規模，因為 $\bar b_x$ 是 $\mu_{x,t}$ 的函數。
這一點解釋了既有文獻的一項觀察：小人口 LC 的問題不因配適期拉長而緩解。

不一致的來源也很清楚：$\Ex[\psi]=0$ 之所以對最小平方自動成立，
前提是損失作用在未經轉換的反應變數上；而式~\eqref{eq:ell} 先取了對數又
替代了零格，兩者都是非線性運算，因而把一個原本自動成立的條件破壞掉了。
把既有文獻以模擬觀察到的偏誤放回 $M$-估計理論中，它的位置就是這裡——
它是任何 $M$-估計量都應該先檢查的第一件事。

既然問題出在母體估計方程不為零，補救的方向也就確定了。
$M$-估計理論對 Fisher 不一致的標準作法是\textbf{中心化估計方程}：
以 $\psi^{*}(u)=\psi(u)-\Ex[\psi(u)]$ 取代 $\psi$，
則 $\Ex[\psi^{*}]=0$ 依定義恢復成立，式~\eqref{eq:pop} 的解也就回到真值。
對平方損失而言 $\psi(u)=u$，故
\begin{equation}\label{eq:center}
\psi^{*}(u_{x,t})=u_{x,t}-b(\mu_{x,t};c_0),
\end{equation}
亦即把每一格的殘差（等價地，每一格的資料）減去 $b$。
第\ref{sec:result}所提的估計量即式~\eqref{eq:center}，它不是另行發明的
校正，而是 $M$-估計中心化在本問題的具體形式。

此處值得指出平方損失在這件事上的一項便利，而這項便利與直覺相反。
因 $\psi(u)=u$ 為線性，把 $u$ 減去一個常數與把 $\psi(u)$ 減去同一個常數
是同一件事，故式~\eqref{eq:center} 的中心化\textbf{精確地}移除了不一致，
無須任何近似。若改用 Huber 損失或檢查函數，$\psi$ 非線性，
$\Ex[\psi(u-b)]\ne\Ex[\psi(u)]-b$，此時要使母體估計方程為零，
所須扣除的量並非 $b$ 本身而是另一個須解方程才能得到的值，
中心化因而只能做到一階。換言之，就「校正 Fisher 不一致」這個目的而言，
\textbf{最不穩健的那個損失反而最好用}——這也是本文雖由穩健統計的框架
出發、最終卻回到平方損失的原因。

由式~\eqref{eq:svd} 與期望值的線性性質，即得 $\hat\alpha_x$ 的偏誤：

\begin{cor}\label{cor:alpha}
$\displaystyle \Ex[\hat\alpha_x]-\alpha_x=\frac{1}{T}\sum_{t=1}^{T} b(\mu_{x,t};c_0).$
\end{cor}

此式與式~\eqref{eq:astar} 形式相同，但兩者的地位不同：式~\eqref{eq:astar}
是漸近的陳述（估計量收斂到哪裡），Corollary~\ref{cor:alpha} 則是\textbf{有限
樣本的恆等式}——只要式~\eqref{eq:pois} 成立、且 $\hat\alpha_x$ 依
式~\eqref{eq:svd} 為列平均，該式對任何 $T$ 與任何 $\mu_{x,t}$ 皆精確成立，
不是漸近近似也不是一階近似。這與小樣本偏誤文獻常見的 $O(n^{-1})$ 型結果
性質不同，原因是 $\hat\alpha_x$ 在奇異值分解之前即已由列平均定出，
不牽涉任何非線性函數的展開；$\hat\beta_x$ 與 $\hat\kappa_t$ 則不然，
兩者是奇異向量，其偏誤須另行處理，見第\ref{sec:betadec}。

\subsection{誤差分布的混合表示}

Proposition~\ref{prop:fisher} 已給出 $b$ 的閉式，但只知道「該扣多少」
還不足以說明這個量在統計上是什麼；而混合模型的角度正好補上這一層，
並且順帶給出兩個可由資料直接計算的門檻。由式~\eqref{eq:ell}，
$\ell_{x,t}$ 的分布為

\begin{equation}\label{eq:mix}
\ell_{x,t}\;\sim\;\pi_{x,t}\,\delta_{\log(c_0/E_{x,t})}
\;+\;(1-\pi_{x,t})\,G_{x,t},
\qquad \pi_{x,t}=e^{-\mu_{x,t}},
\end{equation}
其中 $\delta$ 為點質量、$G_{x,t}$ 為 $\log(D_{x,t}/E_{x,t})$ 在
$D_{x,t}\ge1$ 條件下的分布。亦即\textbf{零格替代使每一格成為一個
兩成分混合}，而式~\eqref{eq:bform} 恰是該混合的期望減去真值：
\begin{equation}\label{eq:bmix}
b(\mu;c_0)=\underbrace{\pi\log c_0}_{\text{退化成分}}
+\underbrace{(1-\pi)\,\Ex[\log D\mid D\ge1]}_{\text{正常成分}}-\log\mu .
\end{equation}
表~\ref{tab:mix} 以數值確認式~\eqref{eq:bmix} 與式~\eqref{eq:bform} 相等
（差異在 $10^{-16}$ 量級，即機器精度），並列出兩個成分的大小。

\begin{table}[H]
\centering
\caption{混合表示的數值確認（$c_0=0.5$）}
\small
\begin{tabular}{rrrrrr}
\toprule
$\mu$ & $\pi=e^{-\mu}$ & $\Ex[\log D\mid D\ge1]$ & 式~\eqref{eq:bform} & 式~\eqref{eq:bmix} & 差\\
\midrule
0.05   & 0.9512 & 0.0174 & $2.33724$  & $2.33724$  & $0$\\
0.30   & 0.7408 & 0.1047 & $0.71762$  & $0.71762$  & $0$\\
0.6931 & \textbf{0.5000} & 0.2433 & $0.14162$  & $0.14162$  & $0$\\
0.9234 & 0.3972 & 0.3245 & $-0.00001$ & $-0.00001$ & $3\times10^{-17}$\\
2.00   & 0.1353 & 0.6934 & $-0.18738$ & $-0.18738$ & $1\times10^{-16}$\\
20.0   & 0.0000 & 2.9696 & $-0.02615$ & $-0.02615$ & $0$\\
\bottomrule
\end{tabular}
\label{tab:mix}
\end{table}

式~\eqref{eq:mix} 有三項後果，而三者都是本文推導的支點。
第一項是它的形式與 Huber 的 $\varepsilon$-污染模型
$(1-\varepsilon)F+\varepsilon H$ 完全相同，但兩者的可知性不同：
古典設定中 $\varepsilon$ 與 $H$ 皆為未知的滋擾，穩健估計量之所以必須
付出效率代價，正是為了在不知道它們的情形下仍提供保障；
而此處 $\pi_{x,t}=e^{-\mu_{x,t}}$ 可由資料估計、$H=\delta$ 完全已知，
既然如此就不必付這筆保險費，而可以直接把污染所造成的後果——
也就是 $b$——算出來扣掉。這是本文選擇偏誤校正而非穩健損失的根本理由，
也是它與第\ref{sec:lit}末一小節所述三項設定的差別所在。
表~\ref{tab:limits} 所列 Huber（1964）一列的「效率代價」限制，
因而在本問題上得以繞過，而非必須承受。

第二項後果是中位數類估計量的失效有解析門檻。
Donoho and Huber（1983）的崩潰點理論指出中位數可承受的污染比例上限為
$50\%$，把它代入式~\eqref{eq:mix} 即得
\begin{equation}\label{eq:ln2}
\pi_{x,t}=e^{-\mu_{x,t}}>\tfrac12
\quad\Longleftrightarrow\quad
\mu_{x,t}<\ln 2=0.6931 .
\end{equation}
亦即以中位數或檢查函數為基礎的估計量，在單年期望死亡數低於 $0.693$ 人的
年齡即失效，而該門檻與 $c_0$ 無關——它是資料的性質，不隨分析者的選擇而動。
此式把一個先前只能以模擬觀察的現象變成可事前計算的量：本文另一份報告記錄
人數五萬時 1 至 4 歲組的零格比例為 $74\%$、已超過崩潰點，而由
式~\eqref{eq:ln2}，該年齡的單年期望死亡數為 $\mu=0.3$，
理論值 $e^{-0.3}=0.7408$，與實測的 $74\%$ 完全相符。
換言之，表~\ref{tab:limits} 中崩潰點一列所載的「$50\%$ 上限」，
在本問題上可以翻譯成一個以死亡數表達的門檻，使用者拿到資料即可判斷
哪些年齡已超出穩健損失的作用範圍。第三項後果則是一個限制而非便利：
式~\eqref{eq:mix} 的第一個成分是點質量而非密度，而有限混合理論
（McLachlan and Peel 2000；Yao and Xiang 2024）通常假設各成分有密度，
這使 EM 演算法無從施加；所幸識別性與標籤交換等標準議題在此同樣不出現，
因為 $\pi$ 由卜瓦松給定、本就不需估計。本文因此不走混合模型的估計路線，
只借用其表示式——這正是表~\ref{tab:limits} 混合模型一列所記的情形。

\subsection{偏誤函數的性質}

式~\eqref{eq:bform} 的形狀完全可由解析得出，而三項性質各對應一個實質的
結論。先看 $\mu$ 趨近於零的一端，此時混合退化為點質量，偏誤為正。
注意 $\log 1=0$ 使式~\eqref{eq:bform} 的 $d=1$ 項消失，故
\begin{equation}\label{eq:small}
b(\mu;c_0)=\log\frac{c_0}{\mu}-\mu\log c_0+O(\mu^2).
\end{equation}
$\mu\to0$ 時 $\pi\to1$，混合退化為點質量，偏誤趨於
$\log(c_0/\mu)\to+\infty$。其意義直接：$\mu<c_0$ 時該格幾乎必然觀測到零，
以 $c_0$ 替代即等於宣稱該格的死亡率為 $c_0/E$，而真值 $\mu/E<c_0/E$，
故必然高估；$c_0=0.5$、$\mu=0.05$ 時式~\eqref{eq:small} 給 $2.33724$，
而式~\eqref{eq:bform} 的精確值為 $2.337237$，兩者差 $6\times10^{-6}$。
這一端所對應的正是圖~\ref{fig:osc}(a) 中那些高懸於真值曲線之上的空心點。

另一端則方向相反。$\mu$ 大時 $\pi\to0$，式~\eqref{eq:bmix} 化為
$\Ex[\log D]-\log\mu$，此時主導的不再是零格替代而是 Jensen 不等式，
偏誤轉為負。
令 $u=(D-\mu)/\mu$，卜瓦松的中心動差給出 $\Ex[u]=0$、$\Ex[u^2]=1/\mu$、
$\Ex[u^3]=1/\mu^2$、$\Ex[u^4]=3/\mu^2+1/\mu^3$，
代入 $\log(1+u)=u-\tfrac{u^2}2+\tfrac{u^3}3-\tfrac{u^4}4+\cdots$ 得
\begin{equation}\label{eq:jensen}
b(\mu;c_0)=-\frac{1}{2\mu}-\frac{5}{12\mu^{2}}+O(\mu^{-3}),
\end{equation}
其中 $-\tfrac5{12}=\tfrac13-\tfrac34$ 由第三、四項合併而來，
且 $c_0$ 在此消失，因為 $\mu$ 大時 $P(D=0)$ 可忽略。
表~\ref{tab:asym} 顯示式~\eqref{eq:jensen} 的精度：二階近似在 $\mu=100$ 時
誤差為 $7.7\times10^{-7}$、在 $\mu=500$ 時為 $6.0\times10^{-9}$，
故對死亡數以百計的高齡組而言，該式已可視為精確。

\begin{table}[H]
\centering
\caption{式~\eqref{eq:jensen} 的精度（$c_0=0.5$）}
\small
\begin{tabular}{rrrrr}
\toprule
$\mu$ & 精確值 & 一階 $-\frac{1}{2\mu}$ & 二階 式~\eqref{eq:jensen} & 二階誤差\\
\midrule
10   & $-0.0552395$ & $-0.0500000$ & $-0.0541667$ & $-1.1\times10^{-3}$\\
20   & $-0.0261518$ & $-0.0250000$ & $-0.0260417$ & $-1.1\times10^{-4}$\\
50   & $-0.0101730$ & $-0.0100000$ & $-0.0101667$ & $-6.4\times10^{-6}$\\
100  & $-0.0050424$ & $-0.0050000$ & $-0.0050417$ & $-7.7\times10^{-7}$\\
500  & $-0.0010017$ & $-0.0010000$ & $-0.0010017$ & $-6.0\times10^{-9}$\\
\bottomrule
\end{tabular}
\label{tab:asym}
\end{table}

兩端的方向既然相反，中間必有一個變號點，而該點是唯一的。
由前兩項性質，$b$ 在 $\mu\to0^{+}$ 時為 $+\infty$、在 $\mu\to\infty$ 時
由負側趨近 $0$；數值上 $b$ 先遞減、於 $\mu=2.2886$ 取極小值 $-0.19104$
後遞增趨近 $0^{-}$，故零點唯一，$c_0=0.5$ 時為 $\mu^\ast=0.9234$。
亦即 $\alpha_x$ 偏誤的符號由該年齡的單年期望死亡數決定：
低於 $0.92$ 人者偏高、高於者偏低。更值得注意的是 $\mu^\ast$ 對 $c_0$ 的
依賴（表~\ref{tab:mustar}）。
偏誤符號的分界線\textbf{不是資料的性質，而是分析者選定替代值的後果}。
$c_0=0.5$ 這個慣例並無統計上的理由，而它決定了哪些年齡被高估。
此點與式~\eqref{eq:ln2} 恰成對比：中位數的失效門檻 $\ln 2$ 是資料的性質，
偏誤的符號門檻 $\mu^\ast$ 則是分析選擇的產物。其後果有二：
不同研究若採不同 $c_0$，所報告的偏誤方向可能相反而不可直接比較；
而 $c_0$ 雖可被選來使某個年齡無偏，卻無法使所有年齡同時無偏——
因 $\mu_{x,t}$ 跨年齡橫跨三個數量級。本文因此不調整 $c_0$ 而直接扣除 $b$。

\begin{table}[H]
\centering
\caption{兩個門檻：一個由資料決定，一個由分析者決定}
\small
\begin{tabular}{llll}
\toprule
門檻 & 意義 & $c_0=0.5$ 時的值 & 是否依賴 $c_0$\\
\midrule
$\mu^\ast$（式~\eqref{eq:bform} 的零點） & $\alpha_x$ 偏誤變號處 & $0.9234$ & \textbf{是}\\
 & \multicolumn{3}{l}{\quad $c_0=0.1/0.3/0.5/1.0\ \Rightarrow\ \mu^\ast=0.134/0.533/0.923/1.509$}\\
$\ln 2$（式~\eqref{eq:ln2}） & 中位數崩潰處 & $0.6931$ & 否\\
\bottomrule
\end{tabular}
\label{tab:mustar}
\end{table}

\subsection{參數偏誤的成分分解}\label{sec:betadec}

Corollary~\ref{cor:alpha} 只涵蓋 $\alpha_x$，因為它在奇異值分解之前即已由
列平均定出。但同一條推導可以再往下走一步，說明 $\beta_x$ 與 $\kappa_t$
受到什麼影響——這一步之所以重要，是因為本文所要回答的問題是「能否降低
$\alpha_x$ 與 $\beta_x$ 的偏誤」，而至此只處理了其中一個。
依式~\eqref{eq:svd}，$(\hat{\boldsymbol\beta},\hat{\boldsymbol\kappa})$ 取自
中心化矩陣 $Z_{x,t}=\ell_{x,t}-\hat\alpha_x$ 的第一組奇異向量，
取其期望並利用 $\sum_t\kappa_t=0$，即得：

\begin{prop}\label{prop:beta}
$\displaystyle \Ex[Z_{x,t}]=\beta_x\kappa_t+\Delta_{x,t},\qquad
\Delta_{x,t}=b(\mu_{x,t};c_0)-\bar b_x,\quad
\bar b_x=\frac1T\sum_s b(\mu_{x,s};c_0).$
\end{prop}

亦即期望矩陣\textbf{不是}秩一訊號 $\boldsymbol\beta\boldsymbol\kappa^{\top}$，
而多了一個確定性的擾動 $\Delta$，其內容為 $b$ 在該年齡\textbf{列內隨 $t$
的變動}。此處有一項值得注意的結構：若 $b(\mu_{x,t})$ 在列內不隨 $t$ 變動，
則 $\Delta=0$，$\beta_x$ 完全不受對數偏誤影響——因為列平均已把它吸收進
$\hat\alpha_x$。但死亡率隨年份改善使 $\mu_{x,t}$ 隨 $t$ 下降，
而 $b$ 是 $\mu$ 的非線性函數，故 $\Delta\ne0$。

由Proposition~\ref{prop:beta} 可把 $\hat\beta$ 的誤差\textbf{精確}分解為兩部分，
不需一階近似：對無噪音的期望矩陣
$\boldsymbol\beta\boldsymbol\kappa^{\top}+\Delta$ 直接作奇異值分解，
所得的 $\hat\beta$ 與真值之差即 $\Delta$ 單獨造成的\textbf{確定性成分}；
以模擬測得的總誤差減去它，即為噪音旋轉奇異向量所造成的\textbf{隨機成分}。
表~\ref{tab:betadec} 給出結果。

\begin{table}[H]
\centering
\caption{$\alpha$、$\beta$ 與漂移項的確定性／隨機成分分解（標準 LC）}
\small
\begin{tabular}{lrrrr}
\toprule
 & $N=10^4$ & $N=5\times10^4$ & $N=2\times10^5$ & $N=10^6$\\
\midrule
$\|\Delta\|_F/\|\boldsymbol\beta\boldsymbol\kappa^{\top}\|_F$
 & 0.693 & 0.499 & 0.189 & 0.064\\
\midrule
$\alpha$ 最大偏誤：確定性 $|\bar b_x|$ & 2.337 & 0.869 & 0.183 & 0.121\\
$\alpha$ 最大偏誤：模擬總量 & 2.331 & 0.895 & 0.296 & 0.212\\
\quad 確定性佔比 & \textbf{100\%} & \textbf{97\%} & 62\% & 57\%\\
\midrule
SSE($\beta$)：確定性 & 0.731 & 0.256 & 0.032 & 0.002\\
SSE($\beta$)：模擬總量 & 7.916 & 4.449 & 1.400 & 0.229\\
\quad 確定性佔比 & \textbf{9.2\%} & \textbf{5.8\%} & 2.3\% & 1.0\%\\
\midrule
漂移偏誤：確定性 & 0.178 & 0.102 & 0.018 & $-0.013$\\
漂移偏誤：模擬總量 & 0.344 & 0.334 & 0.212 & 0.044\\
\quad 確定性佔比 & \textbf{52\%} & 30\% & 8.6\% & ---\\
\bottomrule
\end{tabular}
\label{tab:betadec}
\end{table}

表~\ref{tab:betadec} 是本文對「能不能降低偏誤」這個問題的核心答案，
而三個參數給出的答案並不相同。先看 $\alpha_x$：人數一萬時確定性成分
$2.337$ 與模擬總量 $2.331$ 幾乎相同、佔 $100\%$，人數五萬時為 $97\%$，
這解釋了為何中心化對它有效——所移除的正是全部的偏誤；人數增大後佔比
降至六成，是因為此時絕對偏誤已小，隨機成分的相對重要性隨之上升。
$\beta_x$ 則相反，確定性成分在四個規模下分別只佔總誤差的 $9.2\%$、
$5.8\%$、$2.3\%$、$1.0\%$，其餘九成以上來自異質變異的噪音旋轉主奇異向量，
即各列噪音變異數不同時樣本 Gram 矩陣的對角線被不均勻地灌水，
主向量因而被系統性旋轉。此一機制與 $b$ 無關，因此無論把 $b$ 算得多準
都無法處理它。這是一個否定的結論，但它是有內容的否定：它指出 $\beta_x$
的補救必須改換工具，而正確的工具是加權——若各格的權重使噪音變異數大致
相等，則該旋轉在一階上消失，第\ref{sec:result}的實測證實了這一點。
漂移項介於兩者之間且轉變得很快，人數一萬時確定性成分佔 $52\%$、
五萬時 $30\%$、二十萬時僅 $8.6\%$；故在極小人口下校正對它理論上應有
一半的改善空間，但第\ref{sec:result}將顯示實測並未實現，
因為校正須以 $\hat\mu$ 代入，而 $\hat\mu$ 的誤差所注入的噪音恰好抵銷了
這部分的收益。

表~\ref{tab:betadec} 的數字可以畫成圖，而圖比表更能說明兩個參數的差別
究竟在哪裡。圖~\ref{fig:bband} 把 $\beta_x$ 各估計量的重複分布畫出來，
與圖~\ref{fig:aband} 的 $\alpha_x$ 恰成對照：後者的帶窄而整條偏離真值，
偏誤是主角；前者的帶寬而大致罩住真值，變異數才是主角。
兩張圖並列，等於把表~\ref{tab:betadec} 的確定性佔比由數字翻譯成形狀——
佔比高者表現為「帶窄而偏移」，佔比低者表現為「帶寬而居中」。

\begin{figure}[H]
\centering
\includegraphics[width=\textwidth]{figT2_beta_band.png}
\caption{$\hat\beta_x$ 的 5--95\% 帶（每格一個估計量，上列 $N=10^4$、
下列 $N=5\times10^4$）。黑線為真值。與圖~\ref{fig:aband} 的 $\alpha$ 相反，
$\beta$ 的帶\textbf{很寬}而中線大致罩住真值，是變異數主導的樣子。
中心化單獨反而使帶\textbf{變寬}（$0.743\to0.996$），
須加上資訊加權才收窄（$\to0.233$）。}
\label{fig:bband}
\end{figure}

兩圖並列即本文的核心對照。$\alpha_x$ 的帶窄而偏離真值，
$\beta_x$ 的帶寬而大致罩住真值：人數五萬時 $\alpha$ 的最大偏誤 $0.866$
是其帶寬 $0.218$ 的\textbf{四倍}，而 $\beta$ 的中線離真值中位僅 $0.011$、
帶寬卻達 $0.548$，相差\textbf{五十倍}。這兩個比值正是
表~\ref{tab:betadec} 確定性佔比（$97\%$ 對 $5.8\%$）的另一種表達，
也說明為何同一個校正對兩者的效果截然不同。

圖~\ref{fig:bband} 另顯示兩項須誠實記下的事。其一，中心化單獨使用時
$\beta$ 的帶反而變寬（人數一萬時 $0.743\to0.996$），視覺上比
表~\ref{tab:estcmp} 的 SSE 數字更直接——校正所注入的 $\hat\mu$ 噪音
在 $\beta$ 上是淨損失。其二，中心化加上加權雖使帶寬大幅收窄
（人數五萬時 $0.548\to0.078$，七倍），但其中線在 60 至 85 歲一段
明顯高於真值（中線離真值的中位由 $0.011$ 升至 $0.036$），
亦即它是以\textbf{偏誤換變異數}，而非同時改善兩者。
相較之下 Firth 在人數五萬時的帶寬為 $0.153$、中線離真值僅 $0.003$，
兩者兼得——這也是第\ref{sec:result}中 Firth 在多數指標上仍勝出的原因。

綜合上述，「能不能在某些假設下降低偏誤」這個問題的答案可寫成兩個條件式，
而兩者互不涵蓋。
其一，若\textbf{卜瓦松假設成立}，則 $\alpha_x$ 的偏誤可由式~\eqref{eq:bform}
算出並扣除，改善幅度由表~\ref{tab:betadec} 的確定性佔比給出上界
（人數一萬與五萬時接近全部）。其二，若\textbf{加權後各格的噪音變異數
大致相等}，則 $\beta_x$ 與漂移項的隨機成分可在一階上移除；
但此條件無法由本文的閉式提供，須另行處理。兩個條件互不涵蓋，
這也是小人口 LC 需要兩種藥的原因。

\subsection{與既有修正方法的關係}

最後把本文的校正放回式~\eqref{eq:general} 的分類中，以釐清它與既有補救的
關係。先把該式的損失函數寫清楚。M-分位數的損失為
\begin{equation}\label{eq:mq}
\rho_{\tau,c}(u)=2\,w_\tau(u)\,\psi_c(u),\qquad
w_\tau(u)=\begin{cases}\tau & u>0\\ 1-\tau & u\le 0\end{cases},\qquad
\psi_c(u)=\begin{cases}u^2/2 & |u|\le c\\ c|u|-c^2/2 & |u|>c,\end{cases}
\end{equation}
其中 $\tau$ 控制非對稱、$c$ 控制穩健程度。其兩個極限可直接驗證：
$c\to\infty$ 時 $\psi_c(u)=u^2/2$ 對所有 $u$ 成立，故
$\rho_{\tau,c}(u)=w_\tau(u)u^2$，即 Newey and Powell（1987）的非對稱
最小平方，而 $\tau=0.5$ 時 $w_\tau\equiv0.5$，$\rho$ 退化為 $u^2/2$；
$c\to0$ 時 $\psi_c(u)\to c|u|$，故 $\rho_{\tau,c}(u)\to 2c\,w_\tau(u)|u|$，
而 $w_\tau(u)|u|=u\{\tau-\mathbf{1}(u<0)\}$ 恰為 Koenker and Bassett
（1978）的檢查函數，常數 $2c$ 不影響極小值的位置。
把 $\tau=0.5$、$c=\infty$、$w\equiv1$、$\lambda=0$ 代入式~\eqref{eq:general}
即得式~\eqref{eq:lcobj}，\textbf{標準 LC 因此是該族的一個角落}。

由此角落出發，小人口的各種補救各自改動了其中一項設定：把 $c$ 調成有限或把
$\tau$ 調離 $0.5$ 即得穩健損失；把 $w$ 取為 $\hat\mu$ 即得資訊加權；
把 $\lambda$ 調離零即得 Ridge、LASSO 或經驗貝氏收縮。
問題是這三者與本文的中心化是什麼關係。

關鍵的觀察是三項設定都不改變母體估計方程在真值處的值。就損失形狀而言，
把 $\psi$ 換成 $\psi_{\tau,c}=\rho_{\tau,c}'$ 之後，母體估計方程變為
$\Ex[\psi_{\tau,c}(u(\theta_0))]$，其值一般\textbf{仍不為零}——
事實上因 $\psi_{\tau,c}$ 非線性，該值與 $b$ 的關係還更複雜，
故問題不但沒解決，連「該扣多少」都不再有閉式。就權重而言，
式~\eqref{eq:pop} 在加權後為 $\sum_{x,t}w_{x,t}\Ex[\psi(u_{x,t})]=
\sum_{x,t}w_{x,t}\,b(\mu_{x,t};c_0)$，這是 $b$ 的加權和；
除非權重恰好使該和為零（而那需要事先知道 $b$，即已經做了本文的計算），
否則它同樣不為零。就懲罰而言，$\lambda P(\theta)$ 在解的位置上加一個朝
收縮目標的拉力，這改變的是 $\theta^{\ast}$ 相對於母體估計方程解的位置，
而非母體估計方程本身；它以一個\textbf{新的}偏誤（收縮偏誤）去交換變異數，
並不移除原有的那一項。

三者因此都作用在一個\textbf{已經 Fisher 不一致}的估計方程上，
而式~\eqref{eq:center} 的中心化則直接處理不一致本身。
用式~\eqref{eq:pop} 的語言說，三項設定改變的是估計方程的\textbf{形狀}
（$\psi$）與\textbf{權重}（$w$），中心化改變的是它的\textbf{位置}，
兩者在該式中出現在不同的地方，故彼此正交。

\subsection{加權的最適形式}\label{sec:weight}

式~\eqref{eq:general} 的三項設定中，權重 $w$ 是唯一在本文最終建議中被實際
調動的一個，故其最適形式值得單獨推導。設損失為 $\rho$、$\psi=\rho'$，
逐格權重為 $w_{x,t}$，則加權後的估計方程為
$\sum_{x,t}w_{x,t}\psi(u_{x,t})\,\partial u_{x,t}/\partial\theta=0$。
依 $M$-估計的標準三明治公式，其漸近變異數為 $A^{-1}BA^{-1}$，其中
\[
A=\sum_{x,t} w_{x,t}\,\Ex[\psi'(u_{x,t})]\,g_{x,t}g_{x,t}^{\top},\qquad
B=\sum_{x,t} w_{x,t}^{2}\,\Ex[\psi(u_{x,t})^{2}]\,g_{x,t}g_{x,t}^{\top},
\]
而 $g_{x,t}=\partial u_{x,t}/\partial\theta$。對 $w$ 極小化該變異數，
即得
\begin{equation}\label{eq:wopt}
w\;\propto\;\frac{\Ex[\psi'(u)]}{\Ex[\psi(u)^{2}]} .
\end{equation}
式~\eqref{eq:wopt} 是一個一般結果，把不同的損失代入即得各自的最適權重。
設 $u\sim N(0,\sigma^2)$，三種損失的結果列於表~\ref{tab:wcases}。

\begin{table}[H]
\centering
\caption{三種損失的變異數最適權重
（$\beta_z=\Ex[\min(Z^2,z^2)]$，$Z\sim N(0,1)$；$f$ 為 $u$ 的密度）}
\small
\begin{tabular}{lllll}
\toprule
損失 & $\psi(u)$ & $\Ex[\psi'(u)]$ & $\Ex[\psi(u)^2]$ & 最適權重\\
\midrule
最小平方 & $u$ & $1$ & $\sigma^2$ & $1/\sigma^2$\\
檢查函數 & $\tau-\mathbf{1}(u<0)$ & $f(0)$ & $\tau(1-\tau)$ & $f(0)$\\
Huber$(c)$ & $\min(|u|,c)\,\mathrm{sgn}(u)$ & $2\Phi(c/\sigma)-1$
  & $\sigma^2\beta_{c/\sigma}$ & $\dfrac{2\Phi(c/\sigma)-1}{\sigma^2\beta_{c/\sigma}}$\\
\bottomrule
\end{tabular}
\label{tab:wcases}
\end{table}

第一列即加權最小平方的標準結果。第二列則是 Koenker（2005，第 5.3 節）
定理 5.1 的內容，值得特別指出，因為它與直覺相反：對檢查函數而言，
權重應取該分位處的\textbf{局部密度}而非標準差的倒數，兩者一般並不相同；
本文原先的實作即誤用了後者。第三列的 Huber 損失介於兩者之間，
其形式同時含 $\sigma$ 與截點 $c$，因而是三者中唯一能顯示兩端如何銜接的。
把第三列移到本問題上即得本文所需的形式：由卜瓦松的 delta 方法，
$\mathrm{Var}(\log\hat m_{x,t})\approx1/\mu_{x,t}$，
故 $\sigma_x=\mu_x^{-1/2}$、$c/\sigma_x=c\sqrt{\mu_x}$，代入得
\begin{equation}\label{eq:woptmu}
w_x\;\propto\;\mu_x\,\frac{2\Phi(c\sqrt{\mu_x})-1}{\beta_{c\sqrt{\mu_x}}} .
\end{equation}
式~\eqref{eq:woptmu} 的兩個極限說明了它的結構。$c\sqrt{\mu}\to\infty$ 時
$2\Phi-1\to1$ 且 $\beta\to1$，故 $w\propto\mu$，即最小平方的權重；
$c\sqrt{\mu}\to0$ 時令 $z=c\sqrt{\mu}$，由泰勒展開
$2\Phi(z)-1\approx z\sqrt{2/\pi}$，而 $\beta_z=\Ex[\min(Z^2,z^2)]\approx z^2$
（因 $z$ 小時 $\min(Z^2,z^2)$ 幾乎總是 $z^2$），故
$w\propto\mu\cdot z/z^{2}=\mu/z\propto\sqrt{\mu}$，即檢查函數的權重——
這與第二列一致，因為 $\log\hat m$ 在中位數處的密度為 $\sqrt{\mu/2\pi}$。

由此可得一項結構性的結論：最適權重在 $\sqrt{\mu}$ 與 $\mu$ 之間內插，
而\textbf{內插的變數是 $c\sqrt{\mu}$ 而非 $c$}。因此固定冪次 $w=\mu^{p}$
原則上不可能最適——同一筆資料的各年齡 $\mu$ 橫跨三個數量級，
$c\sqrt{\mu}$ 因而橫跨兩個數量級，式~\eqref{eq:woptmu} 的隱含冪次在
$c=1.345$ 時由 $0.61$ 漂移至 $0.87$，單一冪次只能是某種折衷。
此一預測的實測結果見第\ref{sec:result}末。

\subsection{分布假設的敏感度分析}\label{sec:normal}

上述推導全部建立在式~\eqref{eq:pois} 的卜瓦松假設上，而 $M$-估計的既有
理論多以連續、常態的誤差為設定；既然文獻在常態一側發展得更成熟，
一個自然的問題是能不能改以常態假設推導，以便直接套用那些結果。
答案是一半可以、一半不可以，而分界線恰好落在本問題最關鍵的地方。
可以的那一半是\textbf{加權}。第\ref{sec:weight}的最適權重推導
（式~\eqref{eq:wopt} 與表~\ref{tab:wcases}）全程假設 $u\sim N(0,\sigma^2)$，
其結果與卜瓦松無關，只用到 $\mathrm{Var}(\log\hat m_{x,t})\approx1/\mu_{x,t}$
這個由 delta 方法得來的近似。該部分因而完全落在既有文獻的常態框架內。

不可以的那一半是\textbf{偏誤}。若設
$\ell_{x,t}\sim N\!\big(\log m_{x,t}-\tfrac{1}{2\mu_{x,t}},\,1/\mu_{x,t}\big)$
——這是 $\log(D/E)$ 在 $\mu$ 大時的標準常態近似——則母體估計方程在真值處
的值為
\begin{equation}\label{eq:bnormal}
b_N(\mu)=-\frac{1}{2\mu},
\end{equation}
亦即式~\eqref{eq:jensen} 的一階項。此式簡潔且完全落在常態框架內，
但它\textbf{沒有零格分支}——常態分布取到零的機率為零，
「以 $c_0$ 替代」這件事在常態模型裡根本不存在，因而無從表達。
表~\ref{tab:normal} 把兩者並列。

\begin{table}[H]
\centering
\caption{常態假設下的 $b_N(\mu)=-1/(2\mu)$ 與卜瓦松精確值的比較（$c_0=0.5$）}
\small
\begin{tabular}{rrrrrr}
\toprule
$\mu$ & $P(D=0)$ & 卜瓦松 式~\eqref{eq:bform} & 常態 式~\eqref{eq:bnormal}
 & 差 & 常態的相對誤差\\
\midrule
0.1   & 0.905 & $1.6787$  & $-5.0000$ & $6.679$  & $-398\%$\\
0.3   & 0.741 & $0.7176$  & $-1.6667$ & $2.384$  & $-332\%$\\
0.5   & 0.607 & $0.3416$  & $-1.0000$ & $1.342$  & $-393\%$\\
0.9234& 0.397 & $0.0000$  & $-0.5415$ & $0.542$  & ---\\
2.0   & 0.135 & $-0.1874$ & $-0.2500$ & $0.063$  & $-33\%$\\
5.0   & 0.007 & $-0.1199$ & $-0.1000$ & $-0.020$ & $+17\%$\\
10.0  & 0.000 & $-0.0552$ & $-0.0500$ & $-0.005$ & $+9\%$\\
100.0 & 0.000 & $-0.0050$ & $-0.0050$ & $0.000$  & $+0.8\%$\\
\bottomrule
\end{tabular}
\label{tab:normal}
\end{table}

表~\ref{tab:normal} 的意思很清楚：常態假設給出的偏誤公式在 $\mu>9.5$ 時
相對誤差小於一成，對死亡數以十計以上的年齡完全夠用；
但在 $\mu<1$ 的年齡，它不僅量級錯誤，\textbf{連正負號都相反}——
$\mu=0.3$（即人數五萬時的幼年組）時常態給 $-1.667$ 而實際為 $+0.718$。
原因是常態框架看不見零格替代這個機制，只保留了 Jensen 不等式那一項，
而在小 $\mu$ 處主導的恰是前者。

\textbf{因此本文的取捨是：加權的部分用常態，偏誤的部分必須用卜瓦松。}
這個分工並非權宜，而是反映兩者所刻畫的東西不同——加權處理的是變異數的
異質，那是二階動差的性質，常態近似已足；偏誤處理的是轉換與替代對一階
動差的影響，而替代規則的效果只在離散分布上才存在。
這一點也說明為何既有文獻在常態框架下發展得更成熟卻未觸及本文的問題：
在其標準應用場景中，反應變數本就是連續的，不會出現「必須把零替換掉」
這個步驟。

此一分類給出兩項可檢驗的預測，第\ref{sec:result}將逐一檢驗。
其一，中心化應可大幅移除 $\alpha_x$ 的偏誤，因為 Fisher 不一致正發生在
式~\eqref{eq:svd} 定義 $\hat\alpha_x$ 的那個平均上。
其二，中心化\textbf{不應}改善時間參數的漂移項——漂移項的誤差來自
等權重的秩一分解對異質變異的敏感，那是 $w$ 這個設定的事，
不是中心的事。

\section{以臺灣資料驗證閉式}\label{sec:valid}

Corollary~\ref{cor:alpha} 是一個恆等式，因此嚴格說來並不需要驗證：
只要式~\eqref{eq:pois} 的卜瓦松假設成立、且 $\hat\alpha_x$ 依
式~\eqref{eq:svd} 定義為列平均，該式就必然成立。但這兩個前提在實際的估計
流程中未必如字面那樣乾淨。一方面，實際的 LC 配適含奇異值分解與識別條件的
正規化，而 $\sum_x\hat\beta_x=1$ 與 $\sum_t\hat\kappa_t=0$ 這兩個約束是在
分解之後才施加的，其重新尺度化是否會把 $\hat\alpha_x$ 也帶動，
單看推導並不明顯；另一方面，推導把 $\beta_x\kappa_t$ 當成已知的真值，
而實際上它們也是估計出來的，兩者的誤差是否會滲入 $\hat\alpha_x$，
同樣須以數值確認。本節因此不是在檢驗數學，而是在檢驗\textbf{推導所描述的
那個量，是否確實就是實際估計流程所產生的那個量}。

驗證的資料取自 Human Mortality Database 的臺灣女性 2001 至 2024 年死亡數與
曝露數，聚合為 22 個五齡組，最高年齡組為 100 歲以上。選擇女性是因為其
死亡率曲線較平滑、偏離秩一結構的程度較小，可使驗證聚焦於推導本身而非
模型誤設；選擇 2001 年起則是為避免更早期登記制度變動所帶來的干擾。
模型誤設的影響另於第\ref{sec:real}單獨處理。

真值的定義是本節設計中最須說明的一環。本文以全國資料的卜瓦松配適
$(\alpha^{0},\beta^{0},\kappa^{0})$ 為真值，再按各年齡的曝露結構等比例
縮放至人數 $N\in\{10^4,5\times10^4,2\times10^5\}$，依式~\eqref{eq:pois}
產生小人口的死亡數。此一設計有兩項代價須先承認。其一，它假設小人口的
死亡率\textbf{恰好}等於全國的死亡率，而真實的鄉鎮市區顯然不是如此；
但本節的目的是檢驗閉式，而閉式所刻畫的是估計程序對給定 $\mu_{x,t}$ 的
反應，與該 $\mu_{x,t}$ 是否貼近某個真實地區無關。其二，它使資料恰好落在
秩一雙線性結構上，因而排除了模型誤設；這一點在本節是優點——
若連這樣乾淨的情形都對不上，推導就有問題——但也使本節的結論須由
第\ref{sec:real}補上另一半。

各規模重複 100 次，亂數種子固定為
$\text{seed}=20260926+1000\,i+r$，其中 $i$ 為規模的序號、$r$ 為重複的序號；
第\ref{sec:result}與第\ref{sec:joint}的所有模擬沿用同一組種子，
故各表可逐格對照，不同估計量之間的差異不含蒙地卡羅的雜訊。
每一次重複皆以標準 LC 估計，取各年齡 $\hat\alpha_x-\alpha^{0}_x$ 在 100 次
重複中的中位數作為該年齡的實測偏誤，再與 Corollary~\ref{cor:alpha} 的
預測值 $\bar b_x=T^{-1}\sum_t b(\mu_{x,t};c_0)$ 比較。
比較的方式有二：逐年齡的相關係數，以及以預測值為自變數、實測值為應變數的
迴歸斜率；前者檢驗\textbf{形狀}是否一致，後者檢驗\textbf{量級}是否一致。
兩者都接近 $1$ 才算通過。

\begin{table}[H]
\centering
\caption{閉式預測與模擬實測的對照（標準 LC，重複 100 次）}
\small
\begin{tabular}{lrrr}
\toprule
 & $N=10^4$ & $N=5\times10^4$ & $N=2\times10^5$\\
\midrule
逐年齡的相關係數 & \textbf{1.000} & \textbf{1.000} & \textbf{0.988}\\
迴歸斜率（實測對預測） & 1.001 & 0.990 & 0.962\\
殘差量級 & $\sim$0.01 & $\sim$0.01 & $\sim$0.02\\
\bottomrule
\end{tabular}
\label{tab:valid}
\end{table}

表~\ref{tab:valid} 顯示形狀與量級都通過。逐年齡的相關係數在三個規模下
分別為 $1.000$、$1.000$、$0.988$，迴歸斜率為 $1.001$、$0.990$、$0.962$，
兩者都接近 $1$；殘差的量級約為 $0.01$ 至 $0.02$，而重複 100 次時
$\alpha$ 偏誤的蒙地卡羅誤差本就約有一成，以人數五萬時 5 至 9 歲組的
$0.87$ 計即約 $0.09$，故 $0.01$ 的殘差已在雜訊之內。
人數二十萬時相關係數與斜率略低（$0.988$ 與 $0.962$），
原因是該規模下各年齡的絕對偏誤已降到 $0.1$ 以下，
蒙地卡羅誤差相對於待測的量變大，並非推導本身在該處失準。

表~\ref{tab:validage} 進一步列出逐年齡的對照，而這張表同時展示了
第\ref{sec:normal}前所導出的兩個門檻。先看中位數的崩潰門檻
$\mu<\ln2=0.693$：1 至 4 歲、5 至 9 歲、10 至 14 歲三組的單年期望死亡數
皆為 $0.3$，其零格比例 $\pi=e^{-0.3}=0.74$ 已遠超過 $50\%$，
15 至 19 歲組的 $0.5$ 人對應 $\pi=0.61$ 亦然；
這四組因此落在穩健損失的作用範圍之外，而它們恰好也是偏誤最大的四組
（$0.33$ 至 $0.87$）。再看偏誤的變號門檻 $\mu^\ast=0.9234$：
20 至 24 歲組的期望死亡數為 $0.9$，恰在門檻上，其偏誤因而僅 $+0.04$，
幾乎為零；而零歲組的 $1.2$ 人已越過門檻，偏誤轉為 $-0.07$。
兩個門檻所預測的位置與實測一一對應，且都可由死亡數直接算出，
不需要真值——這正是把閉式寫出來的實用價值所在。

\begin{table}[H]
\centering
\caption{逐年齡對照（$N=5\times10^4$，選列）}
\small
\begin{tabular}{lrrrrr}
\toprule
年齡組 & 單年期望死亡數 & $\pi=e^{-\mu}$ & 閉式預測 & 模擬實測 & 差\\
\midrule
0      & 1.2  & 0.30 & $-0.0753$ & $-0.0692$ & $+0.0061$\\
1--4   & 0.3  & \textbf{0.74} & $0.7283$  & $0.7135$  & $-0.0148$\\
5--9   & 0.3  & \textbf{0.77} & $0.8689$  & $0.8658$  & $-0.0031$\\
10--14 & 0.3  & \textbf{0.77} & $0.8341$  & $0.8262$  & $-0.0079$\\
15--19 & 0.5  & \textbf{0.61} & $0.3298$  & $0.3386$  & $+0.0088$\\
20--24 & 0.9  & 0.41 & $0.0415$  & $0.0529$  & $+0.0114$\\
45--49 & 7.3  & 0.00 & $-0.0793$ & $-0.0831$ & $-0.0038$\\
70--74 & 48.8 & 0.00 & $-0.0108$ & $-0.0135$ & $-0.0027$\\
85--89 & 73.9 & 0.00 & $-0.0069$ & $-0.0061$ & $+0.0008$\\
\bottomrule
\end{tabular}
\label{tab:validage}
\end{table}

圖~\ref{fig:aband} 把同一件事畫出來，而圖比表更能說明偏誤與變異的相對
大小。圖中紅色區域為標準 LC 在 100 次重複中 $\hat\alpha_x-\alpha_x$ 的
$5$ 至 $95$ 百分位帶，紅線為其中線，黑色虛線則是Corollary~\ref{cor:alpha}
的閉式預測 $\bar b_x$。

\begin{figure}[H]
\centering
\includegraphics[width=\textwidth]{figT1_alpha_band.png}
\caption{$\hat\alpha_x$ 的偏誤：5--95\% 帶與閉式預測。黑色虛線為
Corollary~\ref{cor:alpha} 的預測值 $\bar b_x$，與標準 LC 的中線幾乎完全重合
（最大差 $0.043$ 與 $0.019$）。標準 LC 的帶\textbf{窄}而整條\textbf{偏離
零}，是偏誤主導的樣子；中心化（藍）把中線拉回零附近，代價是帶變寬。}
\label{fig:aband}
\end{figure}

圖~\ref{fig:aband} 有三點可讀。第一，\textbf{黑色虛線落在紅線上}：
閉式預測與模擬中線的最大差在兩個規模下分別為 $0.043$ 與 $0.019$，
而偏誤本身在 5 至 9 歲組達 $2.33$ 與 $0.87$。這是Corollary~\ref{cor:alpha}
的視覺版本。第二，\textbf{標準 LC 的帶很窄}——人數五萬時帶寬中位為
$0.218$，而 5 至 9 歲組的偏誤為 $0.866$，偏誤是帶寬的四倍。
亦即該年齡幾乎每一次模擬都落在真值的同一側，這正是「偏誤主導」的
操作型定義，也說明為何重複抽樣無法察覺它。第三，中心化（藍）把中線
拉回零附近，但帶變寬（人數五萬時 $0.218\to0.277$）——這是以 $\hat\mu$
代入所注入的噪音，屬第\ref{sec:result}所述的代價。

綜合表~\ref{tab:valid}、表~\ref{tab:validage} 與圖~\ref{fig:aband}，
本節的結論是推導所描述的量確實就是實際估計流程所產生的量：
奇異值分解與識別條件的正規化並未把 $\hat\alpha_x$ 帶動，
$\hat\beta$ 與 $\hat\kappa$ 的估計誤差也未明顯滲入。
這把先前只能以模擬描述的現象變成一個可計算的量——標準 LC 的 $\alpha_x$
偏誤\textbf{完全}由各格期望死亡數的一個確定性函數決定，不含未知成分。
既然如此，它就可以被算出來並扣掉，而這正是下一節所要做的事。
須補充的是，本節的乾淨只是設計出來的乾淨：資料恰好落在秩一結構上、
分布恰好是卜瓦松。真實資料兩者都不完全成立，
故第\ref{sec:real}將以實際觀測的死亡數重做一次，檢驗結論是否仍然成立。

\section{以中心化修改 LC 的估計}\label{sec:result}

\subsection{修改後的估計流程}

式~\eqref{eq:center} 的中心化在實作上是扣除\textbf{格}而非校正參數，
使校正同時流入三個參數。困難在於 $b$ 依賴未知的 $\mu_{x,t}$，故須迭代：

實作上先以標準 LC 起始，得到一組
$(\hat{\boldsymbol\alpha},\hat{\boldsymbol\beta},\hat{\boldsymbol\kappa})$，
據以算出各格的期望死亡數
$\hat\mu_{x,t}=E_{x,t}\exp(\hat\alpha_x+\hat\beta_x\hat\kappa_t)$；
再以 $\ell_{x,t}-b(\hat\mu_{x,t};c_0)$ 取代原始的 $\ell_{x,t}$，
依式~\eqref{eq:svd} 重新配適，得到新的一組參數後回頭更新 $\hat\mu_{x,t}$，
如此反覆至參數的變動小於容忍值為止，實測收斂約需三至六次。
此一迭代之所以必要，是因為式~\eqref{eq:bform} 依賴未知的 $\mu_{x,t}$，
而後者只能由配適值估計；也正因為如此，$\hat\mu$ 的誤差會隨校正一併注入
每一格，構成第\ref{sec:result}末所述的第一項限制。

與偏誤減少文獻的關係值得一提。Firth（1993）把偏誤減少方法分為
\textbf{事後修正}（先算最大概似估計值再扣除偏誤的估計，
即 Cox and Snell 1968）與\textbf{預先修正}（修改估計方程本身，
即 Firth 的 Jeffreys 懲罰）。本文屬前者，但施加對象不同，
而這個不同是決定性的：Cox--Snell 修正的是最大概似估計量的 $O(n^{-1})$
偏誤，前提是最大概似估計值\textbf{存在且有限}；而小人口 LC 的核心困難
正是零格導致的完全分離（Albert and Anderson 1984），
卜瓦松最大概似的 $\hat\alpha_x$ 可為 $-\infty$（本文實測人數一萬時
5 至 9 歲組為 $-13.03$），故 Cox--Snell 在此無從施加。
式~\eqref{eq:bform} 則完全不涉及最大概似估計量，只用到 $D$ 的邊際分布。

\subsection{實證結果}

以下的比較採與第\ref{sec:valid}完全相同的亂數種子，故各表可逐格對照。
所比較的五個估計量涵蓋表~\ref{tab:limits} 的前三列：標準 LC 為未修正的
基準，卜瓦松最大概似用以呈現事後修正一列所載「須最大概似估計值存在」
這項限制在本問題上的後果，Firth 代表預先修正一路，
中心化與中心化加權則為本文所提的修正。

\begin{table}[H]
\centering
\caption{中心化與既有作法的比較（重複 100 次，取中位數）}
\small
\begin{tabular}{llrrrrr}
\toprule
$N$ & 估計量 & $\alpha$ 中位 & $\alpha$ 最大 & SSE($\beta$) & cor($\beta$) & 漂移偏誤\\
\midrule
\multirow{5}{*}{$10^4$}
 & 標準 LC          & 0.1612 & 2.3309 & 7.916 & $-0.074$ & 0.3436\\
 & 卜瓦松 MLE       & 0.0495 & 13.0317 & 15.525 & 0.118 & 0.1612\\
 & Firth            & \textbf{0.0248} & \textbf{0.4267} & 17.797 & 0.116 & 0.2237\\
 & 中心化           & 0.0693 & 0.5008 & 18.048 & $-0.078$ & 0.3544\\
 & 中心化 $+$ 加權   & 0.1749 & 0.7399 & \textbf{2.537} & 0.111 & \textbf{0.1884}\\
\midrule
\multirow{5}{*}{$5\times10^4$}
 & 標準 LC          & 0.0611 & 0.8658 & 4.449 & 0.037 & 0.3339\\
 & 卜瓦松 MLE       & 0.0073 & 0.2260 & 1.756 & 0.365 & 0.0492\\
 & Firth            & 0.0088 & \textbf{0.0680} & 1.410 & \textbf{0.397} & \textbf{0.0485}\\
 & 中心化           & \textbf{0.0093} & 0.2434 & 8.834 & 0.198 & 0.3335\\
 & 中心化 $+$ 加權   & 0.0110 & 0.2524 & \textbf{1.023} & 0.179 & 0.1381\\
\midrule
\multirow{5}{*}{$2\times10^5$}
 & 標準 LC          & 0.0231 & 0.1846 & 1.400 & 0.379 & 0.2117\\
 & 卜瓦松 MLE       & \textbf{0.0018} & 0.0361 & 0.346 & 0.596 & 0.0267\\
 & Firth            & 0.0025 & \textbf{0.0398} & \textbf{0.333} & \textbf{0.605} & \textbf{0.0262}\\
 & 中心化           & 0.0032 & 0.1778 & 3.499 & 0.438 & 0.2747\\
 & 中心化 $+$ 加權   & 0.0071 & 0.2897 & 0.832 & 0.213 & 0.1106\\
\bottomrule
\end{tabular}
\label{tab:estcmp}
\end{table}

圖~\ref{fig:cmp} 與圖~\ref{fig:e0} 把表~\ref{tab:estcmp} 與
表~\ref{tab:real} 的內容畫出來，前者比較三項參數指標、後者直接看使用者
真正關心的平均餘命。

\begin{figure}[H]
\centering
\includegraphics[width=\textwidth]{figT3_estimator_compare.png}
\caption{三項指標隨人口規模的變化（縱軸取對數）。實線與實心點為卜瓦松
模擬，虛線與空心點為實際觀測死亡數的二項抽薄。兩組幾乎平行，
顯示結論不因本資料所含的模型誤設而改變。}
\label{fig:cmp}
\end{figure}

\begin{figure}[H]
\centering
\includegraphics[width=0.78\textwidth]{figT4_e0_by_estimator.png}
\caption{配適末年零歲平均餘命的 100 次重複（$N=5\times10^4$）。
虛線為真值 $84.07$ 歲，各估計量上方標示其中位偏誤。標準 LC 低估 $1.48$ 歲，
中心化加權則為 $+0.02$ 歲。}
\label{fig:e0}
\end{figure}

圖~\ref{fig:cmp} 有一項只有並列才看得出來的訊息：實線與虛線幾乎平行，
亦即以真實觀測死亡數抽薄所得的結論與純卜瓦松模擬一致；
唯一的例外是標準 LC 的 SSE($\beta$)，其虛線明顯高於實線，
這正是第\ref{sec:real}所述「真實資料中不合秩一結構的成分也被 $\hat\beta$
吸收」的表現。圖~\ref{fig:e0} 則把參數層次的改善翻譯成使用者真正關心的量：
標準 LC 的一百次重複全部落在真值之下，中位數低估 $1.48$ 歲；
卜瓦松最大概似與 Firth 分別改善至 $-0.20$ 與 $-0.28$ 歲；
而中心化加權的中位數為 $+0.02$ 歲，點雲已大致對稱地分布在真值兩側。

第\ref{sec:method}末所提的兩項預測都成立，而其中第二項的成立方式比第一項
更有說服力。第一項預測是中心化應大幅移除 $\alpha_x$ 的偏誤：
$\alpha$ 偏誤中位數在三個規模下分別下降 $2.3$、
$6.6$、$7.2$ 倍；人數一萬時 $\alpha$ 最大偏誤由 $2.3309$ 降至 $0.5008$，
且決定性地優於卜瓦松最大概似的 $13.03$——後者因完全分離而崩壞，
這正是式~\eqref{eq:bform} 不依賴最大概似這個性質的價值。
加上資訊加權之後，該估計量在\textbf{三個規模的每一項指標上都優於標準 LC}
（$\beta$ 的平方誤差和 $7.92\to2.54$、$4.45\to1.02$、$1.40\to0.83$；
漂移項偏誤 $0.34\to0.19$、$0.33\to0.14$、$0.21\to0.11$），
且完全留在奇異值分解的架構內、沒有調校參數。

第二項預測是中心化\textbf{不應}改善漂移項，而這一項的成立是本文分工論
最強的證據。漂移項偏誤由 $0.3436$、$0.3339$、$0.2117$ 變為 $0.3544$、
$0.3335$、$0.2747$，即無改善。這與第\ref{sec:method}末的預測一致：
中心化動的是估計方程的中心，而漂移項的誤差來自權重。
既然一個確定性的 $\mu$ 函數可以移除 $\alpha$ 的偏誤卻完全不動漂移項，
那麼漂移項的偏誤就不可能是對數轉換造成的——它須以 $w$ 這個設定處理，
這也正是「中心化 $+$ 加權」一列改善漂移項的原因。

\subsection{模型誤設下的穩健性檢驗}\label{sec:real}

前一小節的模擬取 $D\sim\text{Poisson}(E\,m^{\text{fitted}})$，
亦即假設 LC 模型完全正確，小人口的資料因而不含任何不合模型的結構。
這對檢驗式~\eqref{eq:bform} 本身是恰當的，但無法回答一個實務上更重要的
問題：若真實資料本來就不完全服從秩一雙線性結構，校正是否仍有效？

為此本文另設計一項以真實資料為基礎的檢驗。取\textbf{實際觀測到的}全國
死亡數 $D^{\text{obs}}_{x,t}$，以二項抽薄產生小人口：
\begin{equation}\label{eq:thin}
D_{x,t}\sim\text{Binomial}\big(D^{\text{obs}}_{x,t},\,p\big),\qquad
E_{x,t}=p\,E^{\text{obs}}_{x,t},\qquad p=N/N^{\text{obs}} .
\end{equation}
如此小人口繼承臺灣死亡率真實的逐年不規則性——所有偏離秩一結構的部分
都被保留下來。實測顯示該偏離並非可忽略：
$\log(D^{\text{obs}}/E^{\text{obs}})$ 對全國 LC 配適值的殘差標準差為
$0.0700$。真值仍取全國資料的 LC 配適。

須說明式~\eqref{eq:thin} 打破的是什麼、保留的是什麼。二項抽薄保持卜瓦松族
（若 $D^{\text{obs}}\sim\text{Poisson}(\mu)$ 則
$D\sim\text{Poisson}(p\mu)$），故式~\eqref{eq:bform} 的形式不變，
只是 $\mu$ 換成 $p\mu$；被打破的不是卜瓦松假設，而是
\textbf{「$\mu_{x,t}$ 恰好落在秩一雙線性結構上」}這個假設。
這正是要檢驗的那一項，因為 $\hat\mu_{x,t}$ 是由 LC 配適算出的，
若真值不在該結構上，$\hat\mu$ 即有系統性誤差，校正也就跟著偏。

\begin{table}[H]
\centering
\caption{真實資料二項抽薄的結果（重複 100 次，取中位數）}
\small
\begin{tabular}{llrrrrr}
\toprule
$N$ & 估計量 & $\alpha$ 中位 & $\alpha$ 最大 & SSE($\beta$) & cor($\beta$) & 漂移偏誤\\
\midrule
\multirow{5}{*}{\shortstack[l]{$10^4$\\ 零格 216/528}}
 & 標準 LC          & 0.1861 & 2.1676 & 6.731 & $-0.264$ & 0.2456\\
 & 卜瓦松 MLE       & 0.0597 & 9.7204 & 14.837 & 0.151 & 0.2758\\
 & Firth            & \textbf{0.0297} & \textbf{0.2047} & 18.551 & 0.113 & 0.2743\\
 & 中心化           & 0.0629 & 0.3271 & 16.292 & $-0.077$ & 0.3488\\
 & 中心化 $+$ 加權   & 0.1258 & 0.5455 & \textbf{3.965} & 0.093 & \textbf{0.2306}\\
\midrule
\multirow{5}{*}{\shortstack[l]{$5\times10^4$\\ 零格 91/528}}
 & 標準 LC          & 0.0791 & 0.7315 & 7.812 & $-0.055$ & 0.3611\\
 & 卜瓦松 MLE       & \textbf{0.0066} & 0.1804 & 1.591 & 0.343 & \textbf{0.0029}\\
 & Firth            & 0.0099 & \textbf{0.0578} & 1.389 & \textbf{0.375} & 0.0164\\
 & 中心化           & 0.0123 & 0.2205 & 6.244 & 0.225 & 0.3156\\
 & 中心化 $+$ 加權   & 0.0544 & 0.2920 & \textbf{1.090} & 0.188 & 0.0985\\
\midrule
\multirow{5}{*}{\shortstack[l]{$2\times10^5$\\ 零格 22/528}}
 & 標準 LC          & 0.0284 & 0.1632 & 1.131 & 0.500 & 0.2316\\
 & 卜瓦松 MLE       & 0.0036 & 0.0301 & 0.259 & \textbf{0.654} & 0.0128\\
 & Firth            & \textbf{0.0034} & \textbf{0.0190} & \textbf{0.254} & 0.650 & \textbf{0.0115}\\
 & 中心化           & 0.0037 & 0.1343 & 2.915 & 0.469 & 0.2759\\
 & 中心化 $+$ 加權   & 0.0335 & 0.3101 & 0.676 & 0.264 & 0.1002\\
\bottomrule
\end{tabular}
\label{tab:real}
\end{table}

三項觀察。第一，\textbf{校正在真實資料上仍然有效}。$\alpha$ 最大偏誤在
三個規模下由 $2.1676$、$0.7315$、$0.1632$ 降至 $0.3271$、$0.2205$、$0.1343$，
改善 $6.6$、$3.3$、$1.2$ 倍。改善幅度與純卜瓦松模擬（$4.7$、$3.6$、$1.0$ 倍）
同一量級，故式~\eqref{eq:bform} 對本資料所含的模型誤設具有一定的穩健性。
理由與前述一致：抽薄不破壞卜瓦松族，而 $\hat\mu$ 只需大致正確即足以
把 $b$ 算到正確的量級——$b$ 在 $\mu$ 上是平滑的，故 $\hat\mu$ 的中等誤差
不會使 $b(\hat\mu)$ 大幅偏離。

第二，\textbf{$\beta_x$ 的結論不變且更明顯}。中心化在三個規模下使
SSE($\beta$) 由 $6.73$、$7.81$、$1.13$ 變為 $16.29$、$6.24$、$2.92$，
即在兩個規模下明顯變差；而中心化加上加權則降至 $3.97$、$1.09$、$0.68$，
三個規模皆優於標準 LC，改善 $1.7$、$7.2$、$1.7$ 倍。
這與第\ref{sec:betadec}表~\ref{tab:betadec} 的分解一致：
$\beta_x$ 的偏誤九成以上是隨機的，須靠加權而非校正。

第三，\textbf{真實資料使標準 LC 的 $\beta$ 更差，而加權的價值因而更高}。
人數五萬時標準 LC 的 SSE($\beta$) 在真實抽薄下為 $7.81$，而在純卜瓦松模擬下
為 $4.45$；差額來自真實資料中不合秩一結構的成分也被 $\hat\beta$ 吸收。
相對地，加權版本在兩種設計下分別為 $1.09$ 與 $1.02$，幾乎不變。
換言之\textbf{加權不只處理異質變異，也吸收了一部分模型誤設}，
這是模擬設計所看不到的。

\subsection{兩參數的聯合改善}\label{sec:joint}

表~\ref{tab:estcmp} 呈現了一個表面上的取捨：逐格校正把 $\alpha$ 修好卻把
$\beta$ 弄差，加上加權後 $\beta$ 修好了但 $\alpha$ 又退了一些。
這是否表示兩者只能擇一？第\ref{sec:betadec}的分解說不是，而且指出了出路。

由 Proposition~\ref{prop:beta}，逐格扣除 $b$ 同時做了兩件事：一是移除
$\alpha_x$ 的偏誤，其量為列平均 $\bar b_x$；二是移除 $\beta_x$ 的確定性
擾動 $\Delta_{x,t}=b(\mu_{x,t})-\bar b_x$。但表~\ref{tab:betadec} 已顯示
第二件事的收益很小——$\Delta$ 只佔 $\beta$ 總誤差的一成以下——
而逐格扣除所注入的 $\hat\mu$ 噪音卻進入每一格，對依賴跨格相對大小的
$\beta$ 傷害甚大。\textbf{既然如此，正確的設計是只做第一件事}：
把 $\hat\alpha_x$ 扣掉 $\bar b_x$，而完全不動中心化矩陣。
由式~\eqref{eq:svd}，$\hat\alpha_x$ 是列平均、在奇異值分解之前即已定出，
故事後扣除 $\bar b_x$ 對 $\hat\beta_x$ 與 $\hat\kappa_t$ 的計算毫無影響；
兩項改善因而可以疊加而不互相干擾。表~\ref{tab:joint} 檢驗此一設計。

\begin{table}[H]
\centering
\caption{逐格校正與僅校正 $\alpha$ 的比較（$N=5\times10^4$，重複 100 次）}
\small
\begin{tabular}{lrrrrr}
\toprule
估計量 & $\alpha$ 中位 & $\alpha$ 最大 & SSE($\beta$) & 漂移偏誤 & $e_0$ 偏誤（歲）\\
\midrule
標準 LC              & 0.0611 & 0.8658 & 4.4494 & 0.3339 & $-1.480$\\
Firth                & 0.0088 & 0.0680 & 1.4102 & \textbf{0.0485} & $-0.282$\\
逐格校正             & 0.0093 & 0.2434 & 8.8335 & 0.3335 & $-1.522$\\
逐格校正 $+$ 加權     & 0.0110 & 0.2524 & \textbf{1.0230} & 0.1381 & $+0.023$\\
僅校正 $\alpha$      & \textbf{0.0063} & \textbf{0.0439} & 4.4494 & 0.3339 & $-1.466$\\
僅校正 $\alpha$ $+$ 加權 & 0.0238 & 0.0715 & 1.1497 & 0.1407 & \textbf{$+0.007$}\\
\bottomrule
\end{tabular}
\label{tab:joint}
\end{table}

結果證實了上述推理，而且有兩處值得特別指出。其一，僅校正 $\alpha$ 的版本
其 SSE($\beta$) 為 $4.4494$，與標準 LC \textbf{完全相同到小數第四位}——
這不是巧合，而是因為該作法根本沒有碰到中心化矩陣，$\hat\beta$ 的計算
逐位元一致。其二，該版本的 $\alpha$ 最大偏誤為 $0.0439$，\textbf{優於
Firth 的 $0.0680$}，是全表最低；相對地逐格校正為 $0.2434$，
可見逐格扣除所注入的噪音同時也傷害了 $\alpha$ 本身。
把它與資訊加權疊加之後，$\alpha$ 最大偏誤 $0.0715$ 與 Firth 相當、
SSE($\beta$) $1.1497$ 亦與 Firth 的 $1.4102$ 相當，
而零歲平均餘命的偏誤為 $+0.007$ 歲，是全表最接近零者。
\textbf{因此兩者可以同時改善，條件是校正只施加在它該施加的地方。}

圖~\ref{fig:pareto} 把這件事畫成一張取捨圖，橫軸為 SSE($\beta$)、
縱軸為 $\alpha$ 最大偏誤，皆取對數，愈靠左下愈好。

\begin{figure}[H]
\centering
\includegraphics[width=\textwidth]{figT6_pareto.png}
\caption{$\alpha$ 與 $\beta$ 的取捨圖。由標準 LC 出發，僅校正 $\alpha$ 是
\textbf{垂直向下}（只改善 $\alpha$，$\beta$ 完全不動）、逐格校正加權是
\textbf{水平向左}（主要改善 $\beta$），而僅校正 $\alpha$ 加上加權則是
\textbf{對角線}（兩者同時）。$N=5\times10^4$ 時該點與 Firth 幾乎重合。}
\label{fig:pareto}
\end{figure}

圖~\ref{fig:pareto} 的三個箭頭是本小節的摘要。垂直的虛線箭頭表示僅校正
$\alpha$ 只在縱軸上移動，橫軸位置分毫未變；水平的虛線箭頭表示加權主要在
橫軸上移動；而實線的對角箭頭表示兩者疊加後同時往左下走。
人數二十萬時 Firth 仍在更左下的位置，但在五萬與一萬時兩者已相當接近，
而後者正是鄉鎮市區的實際規模。

圖~\ref{fig:jointband} 則把同一件事回到逐年齡的層次。

\begin{figure}[H]
\centering
\includegraphics[width=\textwidth]{figT7_joint_band.png}
\caption{標準 LC、Firth 與僅校正 $\alpha$ 加權三者的逐年齡分布
（$N=5\times10^4$，5--95\% 帶）。(a) 綠線在全年齡貼近零，
標準 LC 則在幼年組隆起；(b) 綠色的帶最窄，其中線在多數年齡跟上真值。}
\label{fig:jointband}
\end{figure}

\subsection{三項限制}

本文的修正有三項限制，而三項都可回到表~\ref{tab:limits} 的某一列。
第一項是純中心化使 $\beta$ 變差：表~\ref{tab:estcmp} 中其 SSE($\beta$)
為 $18.05$、$8.83$、$3.50$，均大於標準 LC 的 $7.92$、$4.45$、$1.40$。
機制是逐格扣除 $b(\hat\mu_{x,t})$ 把 $\hat\mu$ 的估計誤差注入每一格，
而 $\beta$ 的估計依賴跨格的相對大小，因而對這種注入特別敏感，
必須配合加權才能抵銷。這是真實的代價而非實作問題，
也正是表~\ref{tab:limits} 中正則化一列所載「以偏誤換變異數」的另一面：
此處換的方向相反，是以變異數換偏誤，而在 $\beta$ 上這筆交易並不划算。
第二項限制是 Firth 在 $\alpha$ 最大偏誤上仍全面勝出，其 $0.4267$、
$0.0680$、$0.0398$ 優於中心化的 $0.5008$、$0.2434$、$0.1778$；
理由與兩者的方法性質一致——中心化只移除期望值的偏移，每一次實現的偏差
仍在，而 Firth 在計數尺度上估計、根本不產生那筆錯誤資訊。
這對應表~\ref{tab:limits} 中預先修正一列的優點，也說明事後與預先兩路
並非等價。本文另掃描了部分扣除（係數 $0.8$ 與 $0.6$），結果一律較完全
扣除為差，故不存在過度校正，原先「$\hat\mu$ 有誤差故應保守扣除」的推測
不成立。第三項限制則與第\ref{sec:weight}的推導有關：該節所導出的最適權重
式~\eqref{eq:woptmu} 雖然正確，實測上卻與簡單的冪次權重沒有可分辨的差異。
以 $\sqrt{\hat\mu}$、$\hat\mu$ 與式~\eqref{eq:woptmu} 三者相比，
$N=2\times10^5$、$c=1.345$ 時 $\beta$ 的平方誤差和分別為 $0.4051$、
$0.4108$、$0.4017$，彼此的差距皆在蒙地卡羅誤差之內；
而同一設定下無加權為 $1.0100$。換言之真正有價值的是「有沒有加權」，
而非「權重的形式對不對」。解釋是目標函數對權重的冪次相當平坦：
權重的作用是反映 $\sigma_x^2$ 跨三個數量級的落差，三者都足以做到這件事，
其間的差異屬二階。實務上因此取最簡單的 $w=\hat\mu$ 即可，
而式~\eqref{eq:woptmu} 的價值在於它說明了為何固定冪次\textbf{原則上}
不可能最適——這一點在推導上重要，在本資料的量級下則不重要。

\section{討論與結論}\label{sec:disc}

本文由 $M$-估計、混合模型與 M-分位數 LASSO 三條一般理論的性質，
推導小人口 LC 模型的偏誤，得到三項結論。

第一，既有文獻以模擬觀察到的 $\alpha_x$ 偏誤，在 $M$-估計的架構中有一個
明確的名字與一個閉式：它是標準 LC 估計方程的\textbf{Fisher 不一致}，
其量為 $b(\mu;c_0)$（Proposition~\ref{prop:fisher}），而 $\hat\alpha_x$ 的偏誤
恰為各格 $b$ 的算術平均（Corollary~\ref{cor:alpha}），且後者是恆等式。
臺灣資料的驗證給出 $1.000$、$1.000$、$0.988$ 的相關係數。
這把「偏誤有多大」由一個須以模擬回答的問題，變成一個可以計算的量。

第二，混合模型的表示式說明了零格替代在統計上是什麼，並給出兩個門檻。
式~\eqref{eq:mix} 顯示每一格是一個混合權重\textbf{已知}的兩成分混合，
因此古典穩健統計為未知污染比例所付的效率代價在此不必付——
這是本文選擇中心化而非穩健損失的理由。而把崩潰點理論代入該表示式，
即得中位數失效的解析門檻 $\mu<\ln2$（式~\eqref{eq:ln2}），
與偏誤變號的門檻 $\mu^\ast=0.9234$ 恰成對比：前者是資料的性質，
後者則隨分析者選定的 $c_0$ 移動。

第三，M-分位數 LASSO 的分類說明中心化與既有補救的關係是正交的，
並給出兩項獲得證實的預測：中心化大幅改善 $\alpha_x$，但不改善
$\beta_x$ 與漂移項。後者本身即是結果——它證明小人口 LC 的偏誤
至少來自兩個獨立機制，$\alpha_x$ 源於估計方程的中心不正，
$\beta_x$ 源於權重不當。

第四，\textbf{就「能不能在某些假設下降低 $\alpha_x$ 與 $\beta_x$ 的小樣本
偏誤」這個問題，本文的答案對兩個參數不同，且差別可以量化。}
第\ref{sec:betadec}把兩者的偏誤精確分解為確定性與隨機成分：
$\alpha_x$ 的偏誤在人數一萬與五萬時有 $100\%$ 與 $97\%$ 是確定性的，
故在卜瓦松假設下可由閉式近乎完全移除；$\beta_x$ 的偏誤則只有
$9.2\%$ 與 $5.8\%$ 是確定性的，其餘來自異質變異噪音旋轉主奇異向量，
\textbf{無論把 $b$ 算得多準都無法處理}，須改以加權為之。
因此可降低偏誤的條件是兩個而非一個：卜瓦松假設可買到 $\alpha_x$，
「加權後噪音變異數大致相等」可買到 $\beta_x$，而後者不由本文的閉式提供。

第五，以真實資料二項抽薄的檢驗顯示上述結論對模型誤設具有穩健性
（第\ref{sec:real}）。$\alpha$ 最大偏誤的改善幅度在真實抽薄與純卜瓦松模擬
下同為 $1.2$ 至 $6.6$ 倍的量級。一項只有真實資料才顯示出來的發現是：
標準 LC 的 SSE($\beta$) 在真實抽薄下較純模擬明顯為差（人數五萬時 $7.81$
對 $4.45$），而加權版本在兩種設計下幾乎不變（$1.09$ 對 $1.02$），
亦即\textbf{加權不只處理異質變異，也吸收了一部分模型誤設}。

把第\ref{sec:limtable}表~\ref{tab:limits} 的各列逐一回顧，可以看出本文的
結果如何落在既有文獻所記載的限制之中。事後修正一列所載「須最大概似估計值
存在」這項限制，在本問題上確實出現且後果嚴重——卜瓦松最大概似在人數一萬時
給出 $-13.03$，而本文由邊際分布推導的作法正是為繞過它而設。
預先修正一列的限制（$O(n^{-1})$ 的漸近性）同樣出現，Firth 的殘餘偏誤
在人數一萬時仍有 $0.427$，只是遠小於未修正的量級。
穩健損失一列的兩項限制則一項被繞過、一項被證實：效率代價因污染比例已知
而不必承受，但「污染為少數」的前提在幼年組明確不成立，
其比例達 $0.74$，超過崩潰點所允許的 $50\%$。
正則化一列的調參限制在本問題上格外嚴重，因為 LC 的跨格借力使交叉驗證
無合法的切分方式；而該列所載「以偏誤換變異數」的取捨則以另一種形式出現
在圖~\ref{fig:bband} 上——加權使 $\beta$ 的帶寬收窄七倍，卻使中線偏離真值。
經驗貝氏收縮一列的「直接估計量不偏」前提亦不完全成立，
因 $\hat\beta_x$ 本身即含 $9.2\%$ 的確定性偏誤。
混合模型一列的限制（成分須有密度）使 EM 無從施加，但因 $\pi$ 已知而
不構成障礙。至於修勻一列，相對於 Yue et al.（2019）以修勻修改 LC 的估計，
本文的修改不引入平滑參數因而不需選參，處理的也不是同一個病灶——
修勻降低的是觀測的變異，本文校正的是估計方程的中心，兩者互補而非替代；
但本文的代價是依賴卜瓦松假設，而修勻並不需要。四個缺口留待後續：
$\hat\beta_x$ 與 $\hat\kappa_t$ 為奇異向量，其偏誤須以矩陣擾動理論處理；
式~\eqref{eq:bform} 依賴卜瓦松假設，實際死亡數若過度離散（例如同一機構
的死亡具相關性），式~\eqref{eq:jensen} 的一階項應為 $-\sigma_D^2/(2\mu^2)$
而非 $-1/(2\mu)$，負二項假設下的對應閉式可同法寫出但尚未做；
本文只處理 $c_0$ 為常數的情形；以及本文以全國配適為真值作模擬，
尚未在真實的鄉鎮市區資料上驗證，而該驗證是本方法能否用於編製鄉鎮市區
生命表的最後一步。

\section*{參考文獻}
\addcontentsline{toc}{section}{參考文獻}

\noindent\textbf{一、中文部分}
\begin{itemize}[leftmargin=2em, label={}, itemsep=3pt]
  \item[] 王信忠、金碩、余清祥（2012）。小區域死亡率推估之研究。人口學刊，
        45：121-154。
  \item[] 王信忠、余清祥、王子瑜（2017）。臺灣原住民死亡率暨生命表編撰研究。
        人口學刊，55：99-131。
  \item[] 余清祥、王信忠、陳譽騰（2021）。年輪變動比用於小區域人口推估的探討。
        人口學刊，63：99-133。
  \item[] 余清祥、王信忠、呂靖翎（2025）。平均餘命與標準化死亡率之相關分析。
        人口學刊，85-116。
  \item[] 陳政勳、余清祥（2010）。小區域人口推估研究：臺北市、雲嘉兩縣、
        澎湖縣的實證分析。人口學刊，41：153-183。
\end{itemize}

\vspace{6pt}
\noindent\textbf{二、英文部分}
\begin{itemize}[leftmargin=2em, itemsep=2pt, parsep=0pt]
  \item[] Bartlett, M. S. 1953. ``Approximate Confidence Intervals.''
        \textit{Biometrika} 40(1-2): 12-19.
  \item[] Albert, A., and Anderson, J. A. 1984. ``On the Existence of Maximum
        Likelihood Estimates in Logistic Regression Models.''
        \textit{Biometrika} 71(1): 1-10.
  \item[] Basellini, U., C. G. Camarda, and H. Booth. 2023. ``Thirty Years on:
        A Review of the Lee-Carter Method for Forecasting Mortality.''
        \textit{International Journal of Forecasting} 39(3): 1033-1049.
  \item[] Bianchi, A., E. Fabrizi, N. Salvati, and N. Tzavidis. 2018.
        ``Estimation and Testing in M-quantile Regression with Applications to
        Small Area Estimation.'' \textit{International Statistical Review}
        86(3): 541-570.
  \item[] Breckling, J., and R. Chambers. 1988. ``M-quantiles.''
        \textit{Biometrika} 75(4): 761-771.
  \item[] Brouhns, N., M. Denuit, and J. K. Vermunt. 2002. ``A Poisson
        Log-bilinear Regression Approach to the Construction of Projected
        Lifetables.'' \textit{Insurance: Mathematics and Economics} 31(3): 373-393.
  \item[] Camarda, C. G., and U. Basellini. 2021. ``Smoothing, Decomposing and
        Forecasting Mortality Rates.'' \textit{European Journal of Population}
        37: 569-602.
  \item[] Chambers, R., and N. Tzavidis. 2006. ``M-quantile Models for Small
        Area Estimation.'' \textit{Biometrika} 93(2): 255-268.
  \item[] Chen, L., A. J. G. Cairns, and T. Kleinow. 2017. ``Small Population
        Bias and Sampling Effects in Stochastic Mortality Modelling.''
        \textit{European Actuarial Journal} 7(1): 193-230.
  \item[] Cox, D. R., and Snell, E. J. 1968. ``A General Definition of Residuals.''
        \textit{Journal of the Royal Statistical Society, Series B} 30(2): 248-265.
  \item[] Currie, I. D., M. Durban, and P. H. C. Eilers. 2004. ``Smoothing and
        Forecasting Mortality Rates.'' \textit{Statistical Modelling} 4(4): 279-298.
  \item[] Dawber, J., N. Salvati, E. Fabrizi, and N. Tzavidis. 2022. ``Expectile
        Regression for Multi-category Outcomes with Application to Small Area
        Estimation of Labour Force Participation.''
        \textit{Journal of the Royal Statistical Society, Series A}.
        doi:10.1111/rssa.12953
  \item[] De Nicol\`o, S., M. R. Ferrante, and S. Pacei. 2023. ``Small Area
        Estimation of Inequality Measures Using Mixtures of Beta.''
        \textit{Journal of the Royal Statistical Society, Series A}.
        doi:10.1093/jrsssa/qnad083
  \item[] Del Sarto, S., M. F. Marino, M. G. Ranalli, and N. Salvati. 2019.
        ``Using Finite Mixtures of M-quantile Regression Models to Handle
        Unobserved Heterogeneity in Assessing the Effect of Meteorology and
        Traffic on Air Quality.'' \textit{Stochastic Environmental Research and
        Risk Assessment} 33: 1345-1359.
  \item[] Dempster, A. P., N. M. Laird, and D. B. Rubin. 1977. ``Maximum
        Likelihood from Incomplete Data via the EM Algorithm.''
        \textit{Journal of the Royal Statistical Society, Series B} 39(1): 1-38.
  \item[] Donoho, D. L., and Huber, P. J. 1983. ``The Notion of Breakdown Point.''
        In \textit{A Festschrift for Erich L. Lehmann}, 157-184. Belmont: Wadsworth.
  \item[] Eilers, P. H. C., and B. D. Marx. 1996. ``Flexible Smoothing with
        B-splines and Penalties.'' \textit{Statistical Science} 11(2): 89-121.
  \item[] Fay, R. E., and Herriot, R. A. 1979. ``Estimates of Income for Small
        Places: An Application of James-Stein Procedures to Census Data.''
        \textit{Journal of the American Statistical Association} 74(366): 269-277.
  \item[] Firth, D. 1993. ``Bias Reduction of Maximum Likelihood Estimates.''
        \textit{Biometrika} 80(1): 27-38.
  \item[] Hampel, F. R. 1974. ``The Influence Curve and Its Role in Robust
        Estimation.'' \textit{Journal of the American Statistical Association}
        69(346): 383-393.
  \item[] Hoerl, A. E., and R. W. Kennard. 1970. ``Ridge Regression: Biased
        Estimation for Nonorthogonal Problems.'' \textit{Technometrics} 12(1): 55-67.
  \item[] Hu, J., Y. Chen, W. Zhang, and X. Guo. 2021. ``Penalized
        High-dimensional M-quantile Regression: From L1 to Lp Optimization.''
        \textit{The Canadian Journal of Statistics} 49(3): 875-905.
  \item[] Huber, P. J. 1964. ``Robust Estimation of a Location Parameter.''
        \textit{Annals of Mathematical Statistics} 35(1): 73-101.
  \item[] Huber, P. J., and Ronchetti, E. M. 2009. \textit{Robust Statistics},
        2nd ed. Hoboken: Wiley.
  \item[] James, W., and C. Stein. 1961. ``Estimation with Quadratic Loss.''
        \textit{Proceedings of the Fourth Berkeley Symposium on Mathematical
        Statistics and Probability} 1: 361-379.
  \item[] Jarner, S. F., and E. M. Kryger. 2011. ``Modelling Adult Mortality in
        Small Populations: The SAINT Model.''
        \textit{ASTIN Bulletin} 41(2): 377-418.
  \item[] Lahiri, P., and N. Salvati. 2023. ``A Nested Error Regression Model
        with High-dimensional Parameter for Small Area Estimation.''
        \textit{Journal of the Royal Statistical Society, Series B} 85(1): 212-239.
  \item[] Lambert, D. 1992. ``Zero-inflated Poisson Regression, with an
        Application to Defects in Manufacturing.''
        \textit{Technometrics} 34(1): 1-14.
  \item[] Liang, K.-Y., and S. L. Zeger. 1986. ``Longitudinal Data Analysis
        Using Generalized Linear Models.'' \textit{Biometrika} 73(1): 13-22.
  \item[] Koenker, R. 2005. \textit{Quantile Regression}. Cambridge: Cambridge
        University Press.
  \item[] Koenker, R., and Bassett, G. 1978. ``Regression Quantiles.''
        \textit{Econometrica} 46(1): 33-50.
  \item[] Lee, R. D., and L. R. Carter. 1992. ``Modeling and Forecasting U.S.
        Mortality.'' \textit{Journal of the American Statistical Association}
        87(419): 659-671.
  \item[] Li, N., and R. Lee. 2005. ``Coherent Mortality Forecasts for a Group of
        Populations: An Extension of the Lee-Carter Method.''
        \textit{Demography} 42(3): 575-594.
  \item[] McLachlan, G., and D. Peel. 2000. \textit{Finite Mixture Models}.
        New York: Wiley.
  \item[] Merlo, L., L. Petrella, N. Salvati, and N. Tzavidis. 2022. ``Marginal
        M-quantile Regression for Multivariate Dependent Data.''
        \textit{Computational Statistics \& Data Analysis} 173: 107500.
  \item[] Newey, W. K., and Powell, J. L. 1987. ``Asymmetric Least Squares
        Estimation and Testing.'' \textit{Econometrica} 55(4): 819-847.
  \item[] Pearson, K. 1894. ``Contributions to the Mathematical Theory of
        Evolution.'' \textit{Philosophical Transactions of the Royal Society of
        London A} 185: 71-110.
  \item[] Ranalli, M. G., N. Salvati, L. Petrella, and F. Pantalone. 2023.
        ``M-quantile Regression Shrinkage and Selection via the Lasso and
        Elastic Net to Assess the Effect of Meteorology and Traffic on Air
        Quality.'' \textit{Biometrical Journal}. doi:10.1002/bimj.202100355
  \item[] Renshaw, A. E., and S. Haberman. 2003. ``Lee--Carter Mortality
        Forecasting with Age-specific Enhancement.''
        \textit{Insurance: Mathematics and Economics} 33(2): 255-272.
  \item[] Bugallo, M., N. Salvati, and D. Morales. 2026. ``Bias Adjustment for
        Mean Squared Error Estimation in M-Quantile Models for Small Area
        Estimation.'' \textit{International Statistical Review}. Forthcoming.
  \item[] Stein, C. 1956. ``Inadmissibility of the Usual Estimator for the Mean
        of a Multivariate Normal Distribution.'' \textit{Proceedings of the
        Third Berkeley Symposium} 1: 197-206.
  \item[] Tibshirani, R. 1996. ``Regression Shrinkage and Selection via the
        Lasso.'' \textit{Journal of the Royal Statistical Society, Series B}
        58(1): 267-288.
  \item[] Wang, H. C., C. J. Yue, and C. T. Chong. 2018. ``Mortality Models and
        Longevity Risk for Small Populations.''
        \textit{Insurance: Mathematics and Economics} 78: 351-359.
  \item[] Wei, Y., A. Pere, R. Koenker, and X. He. 2006. ``Quantile Regression
        Methods for Reference Growth Charts.'' \textit{Statistics in Medicine}
        25(8): 1369-1382.
  \item[] Yao, W., and S. Xiang. 2024. \textit{Mixture Models: Parametric,
        Semiparametric, and New Directions}. Boca Raton: CRC Press.
  \item[] Yue, J. C., H. C. Wang, and T. Y. Wang. 2019. ``Using Graduation to
        Modify the Estimation of Lee-Carter Model for Small Populations.''
        \textit{North American Actuarial Journal}.
        doi:10.1080/10920277.2019.1650288
  \item[] Zoubir, A. M., V. Koivunen, E. Ollila, and M. Muma. 2018.
        \textit{Robust Statistics for Signal Processing}. Cambridge: Cambridge
        University Press.
\end{itemize}

\end{document}
